\chapter{Tipi di dato}
\label{chap:tipi}

Nel capitolo~\ref{chap:calcolatore} abbiamo visto che la memoria contiene
soltanto bit, e che lo stesso byte può essere letto come un numero, come una
lettera o come un pezzo di qualcos'altro: è il \textbf{tipo} a stabilire
quanti byte leggere e come interpretarli. Questo capitolo guarda da vicino i
tipi fondamentali del C: i caratteri, gli interi e i reali. Per ciascuno
vedremo quanta memoria occupa, quali valori può contenere e che cosa succede
quando si superano i suoi limiti; poi come si leggono e si scrivono i
caratteri uno alla volta, come si usano le funzioni matematiche e con quali
regole il C converte un tipo in un altro.

\section{Dichiarazioni, espressioni, assegnazioni}

Tre concetti del capitolo~\ref{chap:primi-programmi} torneranno di
continuo, quindi conviene richiamarli (figura~\ref{fig:dich-espr}):

\begin{itemize}
  \item una \textbf{dichiarazione}, come \code{int i;}, riserva in memoria lo
        spazio per una variabile del tipo indicato e le dà un nome; finché
        non le si assegna nulla, la variabile contiene un valore a caso;
  \item un'\textbf{assegnazione}, come \code{i = 7;}, scrive un valore nella
        variabile;
  \item un'\textbf{espressione}, come \code{i * 2 + 1}, viene calcolata e
        produce un risultato che ha un \textbf{valore} (qui 15) e un
        \textbf{tipo} (qui \code{int}). Anche un'assegnazione è
        un'espressione: \code{i = 7} vale 7.
\end{itemize}

\begin{figure}[H]
  \centering
  \begin{tikzpicture}[font=\Lato\scriptsize,
      tit/.style={font=\Lato\scriptsize\bfseries, text=primary!60!black}]
    \node[tit] at (0.8,1.3) {dichiarazione};
    \node[mcod] at (0.8,0.8) {int i;};
    \variabile{p1}{(1.0,0)}{i}{\mvuota}
    \node[mnota, align=center, anchor=north] at (0.8,-0.4) {si riserva lo spazio;\\il valore è casuale};
    \draw[wtarrow] (2.35,0) -- (3.35,0);
    \node[tit] at (4.9,1.3) {assegnazione};
    \node[mcod] at (4.9,0.8) {i = 7;};
    \variabile{p2}{(5.1,0)}{i}{7}
    \node[mnota, align=center, anchor=north] at (4.9,-0.4) {il valore 7 viene\\scritto nella casella};
    \draw[wtarrow] (6.45,0) -- (7.45,0);
    \node[tit] at (9.4,1.3) {espressione};
    \node[mcod] at (9.4,0.8) {i * 2 + 1};
    \node[wtnodeplain, font=\Lato\scriptsize] at (9.4,0) {valore \textbf{15}\qquad tipo \textbf{int}};
    \node[mnota, align=center, anchor=north] at (9.4,-0.4) {si calcola: il risultato ha\\un valore e un tipo};
  \end{tikzpicture}
  \caption{Dichiarazione, assegnazione ed espressione.}
  \label{fig:dich-espr}
\end{figure}

\section{I tipi fondamentali}

\subsection{Forma completa e forma abbreviata}

Il C ha dodici tipi \term{fondamentali}. Ciascuno ha una forma completa, che
dice tutto esplicitamente, e una forma abbreviata, quella che si usa nella
pratica:

\begin{center}
\renewcommand{\arraystretch}{1.12}
\begin{tabular}{@{}>{\FiraCodeNL}l>{\FiraCodeNL}l@{\hspace{1.6cm}}>{\FiraCodeNL}l>{\FiraCodeNL}l@{}}
  \toprule
  \normalfont\textbf{forma completa} & \normalfont\textbf{si scrive} &
  \normalfont\textbf{forma completa} & \normalfont\textbf{si scrive} \\
  \midrule
  char               & char          & unsigned short int & unsigned short \\
  signed char        & signed char   & unsigned int       & unsigned \\
  unsigned char      & unsigned char & unsigned long int  & unsigned long \\
  signed short int   & short         & float              & float \\
  signed int         & int           & double             & double \\
  signed long int    & long          & long double        & long double \\
  \bottomrule
\end{tabular}
\end{center}

\begin{itemize}
  \item Sono tutte \textbf{parole chiave} (capitolo~\ref{chap:sintassi}):
        non si possono usare come nomi di variabili.
  \item Per \code{short}, \code{int} e \code{long} il segno è sottinteso:
        \code{signed int} è la stessa cosa di \code{int}, quindi la parola
        \code{signed} di solito non si scrive. Anche \code{signed} da solo
        equivale a \code{int}, ma si usa raramente.
  \item \code{short int}, \code{long int} e \code{unsigned int} si
        abbreviano in \code{short}, \code{long} e \code{unsigned}.
  \item Per i caratteri, invece, \code{char}, \code{signed char} e
        \code{unsigned char} sono tre tipi distinti: vedremo fra poco perché.
\end{itemize}

\subsection{Tipi interi e tipi reali}

In base a ciò che rappresentano, i tipi fondamentali formano due famiglie
(figura~\ref{fig:famiglie}). I \term{tipi interi} (\textit{integral types})
contengono numeri interi e, come vedremo, caratteri; i \term{tipi reali}
(\textit{floating types}) contengono numeri con la virgola. Insieme formano i
\term{tipi aritmetici}, quelli su cui si fanno i calcoli.

\begin{figure}[H]
  \centering
  \begin{tikzpicture}[font=\FiraCodeNL\scriptsize,
      tipo/.style={draw=primary!55!black, fill=white, rounded corners=2pt,
        line width=0.6pt, minimum width=2.3cm, minimum height=0.5cm, inner sep=2pt},
      lab/.style={font=\Lato\scriptsize, anchor=west, inner sep=1pt}]
    \draw[rounded corners=6pt, line width=0.8pt, draw=neutral!45!white, fill=neutral!5!white]
      (-1.7,1.55) rectangle (10.25,-2.15);
    \draw[rounded corners=5pt, line width=0.8pt, draw=primary!45!white, fill=primary!7!white]
      (-1.45,0.95) rectangle (6.55,-1.9);
    \draw[rounded corners=5pt, line width=0.8pt, draw=accent!45!white, fill=accent!7!white]
      (7.15,0.95) rectangle (10.0,-1.9);
    \node[lab] at (-1.55,1.28) {\textbf{tipi aritmetici}\ \ {\color{neutral!45!black}ci si fanno i calcoli}};
    \node[lab] at (-1.3,0.67) {\textbf{tipi interi}\ \ {\color{neutral!45!black}numeri interi e caratteri}};
    \node[lab] at (7.3,0.67) {\textbf{tipi reali}};
    \foreach \testo [count=\k from 0] in {char, signed char, unsigned char,
                                          short, int, long,
                                          unsigned short, unsigned, unsigned long} {
      \pgfmathtruncatemacro{\ifam}{\k/3}
      \pgfmathtruncatemacro{\jfam}{mod(\k,3)}
      \node[tipo] at (\jfam*2.55, -\ifam*0.7) {\testo};
    }
    \foreach \testo [count=\k from 0] in {float, double, long double}
      \node[tipo, draw=accent!60!black] at (8.575, -\k*0.7) {\testo};
  \end{tikzpicture}
  \caption{I tipi fondamentali raggruppati per famiglia.}
  \label{fig:famiglie}
\end{figure}

\subsection{I tipi derivati}

A partire dai tipi fondamentali se ne costruiscono altri, i \term{tipi
derivati} (figura~\ref{fig:derivati}): gli \textbf{array} e i
\textbf{puntatori} (capitolo~\ref{chap:puntatori}), le \textbf{strutture}
(capitolo~\ref{chap:strutture}), le enumerazioni e le unioni. In questo
capitolo ci occupiamo solo dei fondamentali.

\begin{figure}[H]
  \centering
  \begin{tikzpicture}[font=\Lato\scriptsize,
      cat/.style={wtnode, font=\Lato\scriptsize\bfseries, inner sep=4pt},
      foglia/.style={font=\Lato\scriptsize, inner sep=2pt, anchor=west}]
    \node[cat] (t) at (0,0) {tipi di dato};
    \node[cat] (f) at (-3.0,-1.15) {fondamentali};
    \node[cat, draw=accent!60!black, fill=accent!8!white] (d) at (3.0,-1.15) {derivati};
    \draw[wtarrow] (t.south) -- ++(0,-0.3) -| (f.north);
    \draw[wtarrow] (t.south) -- ++(0,-0.3) -| (d.north);
    \foreach \testo [count=\k] in {\code{char}, {\code{int} (e \code{short}, \code{long}, \code{unsigned})},
                                   \code{float}, {\code{double} (e \code{long double})}} {
      \node[foglia] (f\k) at ([xshift=16pt, yshift=-\k*0.45cm]f.south west) {\testo};
      \draw[fsramo] ([xshift=8pt]f.south west) |- (f\k.west);
    }
    \foreach \testo [count=\k] in {array (capitolo~\ref{chap:puntatori}),
                                   puntatori (capitolo~\ref{chap:puntatori}),
                                   {strutture, \code{struct} (capitolo~\ref{chap:strutture})},
                                   {enumerazioni, \code{enum}},
                                   {unioni, \code{union}}} {
      \node[foglia] (d\k) at ([xshift=16pt, yshift=-\k*0.45cm]d.south west) {\testo};
      \draw[fsramo] ([xshift=8pt]d.south west) |- (d\k.west);
    }
  \end{tikzpicture}
  \caption{Tipi fondamentali e tipi derivati.}
  \label{fig:derivati}
\end{figure}

\section{Caratteri e il tipo \code{char}}

Cominciamo dal tipo più piccolo, \code{char}, che serve a memorizzare i
caratteri: lettere, cifre, simboli.

\subsection{Un carattere è un numero}

Un computer non conosce le lettere: conosce solo numeri. Per memorizzare un
carattere lo si sostituisce con un numero, secondo una tabella stabilita una
volta per tutte, la \term{codifica}. Un \code{char} occupa \textbf{1 byte},
quindi può assumere $2^8 = 256$ valori diversi: abbastanza per le lettere
maiuscole e minuscole, le cifre, la punteggiatura e alcuni caratteri
speciali come lo spazio, la tabulazione e l'a capo. Non tutti i 256 valori,
però, corrispondono a caratteri stampabili.

In C ogni tipo intero può rappresentare un carattere; in pratica si usano
\code{char} e \code{int}. Le funzioni \code{getchar} e \code{putchar}, che
leggono e scrivono un carattere, lavorano con degli \code{int}; e una
costante carattere come \code{'a'} è di tipo \code{int}, non \code{char}
(in C++ invece è un \code{char}). Lo verificheremo con l'operatore
\code{sizeof}.

\subsection{La codifica ASCII}

Quasi tutti i computer usano la codifica \term{ASCII} (\textit{American
Standard Code for Information Interchange}), che assegna un numero da 0 a 127
a ogni carattere (figura~\ref{fig:ascii}). Esistono altre codifiche, come la
EBCDIC di IBM, usata sui grandi calcolatori (\textit{mainframe}), ma sono rare.

% Una casella della tabella ASCII: \asc{colonna}{riga}{categoria}{carattere}
\newcommand{\asc}[4]{%
  \node[acella, #3] at (#1*2.2, -#2*0.4) {};
  \pgfmathtruncatemacro{\ascn}{#1*16+#2}%
  \node[anum] at (#1*2.2-0.22, -#2*0.4) {\ascn};
  \node[acar] at (#1*2.2+0.42, -#2*0.4) {#4};}
\newcommand{\ascesc}[1]{{\color{accent!75!black}\textbackslash #1}}
\newcommand{\ascnome}[1]{{\fontsize{5}{6}\selectfont\color{neutral!45!black}#1}}

\begin{figure}[H]
  \centering
  \begin{tikzpicture}[
      acella/.style={minimum width=2.1cm, minimum height=0.37cm, rounded corners=1.5pt,
        line width=0.4pt, inner sep=0pt},
      actrl/.style={fill=neutral!9!white, draw=neutral!30!white},
      adig/.style={fill=green-gradient-1!50!white, draw=green-gradient-6!35!white},
      amai/.style={fill=primary!12!white, draw=primary!30!white},
      amin/.style={fill=accent!12!white, draw=accent!30!white},
      apun/.style={fill=white, draw=neutral!30!white},
      anum/.style={font=\FiraCodeNL\tiny, text=neutral!45!black, anchor=east, inner sep=0pt},
      acar/.style={font=\FiraCodeNL\scriptsize\bfseries, inner sep=0pt}]
    \asc{0}{0}{actrl}{\ascesc{0}} \asc{0}{1}{actrl}{\ascnome{SOH}} \asc{0}{2}{actrl}{\ascnome{STX}} \asc{0}{3}{actrl}{\ascnome{ETX}}
    \asc{0}{4}{actrl}{\ascnome{EOT}} \asc{0}{5}{actrl}{\ascnome{ENQ}} \asc{0}{6}{actrl}{\ascnome{ACK}} \asc{0}{7}{actrl}{\ascesc{a}}
    \asc{0}{8}{actrl}{\ascesc{b}} \asc{0}{9}{actrl}{\ascesc{t}} \asc{0}{10}{actrl}{\ascesc{n}} \asc{0}{11}{actrl}{\ascesc{v}}
    \asc{0}{12}{actrl}{\ascesc{f}} \asc{0}{13}{actrl}{\ascesc{r}} \asc{0}{14}{actrl}{\ascnome{SO}} \asc{0}{15}{actrl}{\ascnome{SI}}
    \asc{1}{0}{actrl}{\ascnome{DLE}} \asc{1}{1}{actrl}{\ascnome{DC1}} \asc{1}{2}{actrl}{\ascnome{DC2}} \asc{1}{3}{actrl}{\ascnome{DC3}}
    \asc{1}{4}{actrl}{\ascnome{DC4}} \asc{1}{5}{actrl}{\ascnome{NAK}} \asc{1}{6}{actrl}{\ascnome{SYN}} \asc{1}{7}{actrl}{\ascnome{ETB}}
    \asc{1}{8}{actrl}{\ascnome{CAN}} \asc{1}{9}{actrl}{\ascnome{EM}} \asc{1}{10}{actrl}{\ascnome{SUB}} \asc{1}{11}{actrl}{\ascnome{ESC}}
    \asc{1}{12}{actrl}{\ascnome{FS}} \asc{1}{13}{actrl}{\ascnome{GS}} \asc{1}{14}{actrl}{\ascnome{RS}} \asc{1}{15}{actrl}{\ascnome{US}}
    \asc{2}{0}{apun}{\char"2423} \asc{2}{1}{apun}{!} \asc{2}{2}{apun}{"} \asc{2}{3}{apun}{\#}
    \asc{2}{4}{apun}{\$} \asc{2}{5}{apun}{\%} \asc{2}{6}{apun}{\&} \asc{2}{7}{apun}{'}
    \asc{2}{8}{apun}{(} \asc{2}{9}{apun}{)} \asc{2}{10}{apun}{*} \asc{2}{11}{apun}{+}
    \asc{2}{12}{apun}{,} \asc{2}{13}{apun}{-} \asc{2}{14}{apun}{.} \asc{2}{15}{apun}{/}
    \asc{3}{0}{adig}{0} \asc{3}{1}{adig}{1} \asc{3}{2}{adig}{2} \asc{3}{3}{adig}{3}
    \asc{3}{4}{adig}{4} \asc{3}{5}{adig}{5} \asc{3}{6}{adig}{6} \asc{3}{7}{adig}{7}
    \asc{3}{8}{adig}{8} \asc{3}{9}{adig}{9} \asc{3}{10}{apun}{:} \asc{3}{11}{apun}{;}
    \asc{3}{12}{apun}{<} \asc{3}{13}{apun}{=} \asc{3}{14}{apun}{>} \asc{3}{15}{apun}{?}
    \asc{4}{0}{apun}{@} \asc{4}{1}{amai}{A} \asc{4}{2}{amai}{B} \asc{4}{3}{amai}{C}
    \asc{4}{4}{amai}{D} \asc{4}{5}{amai}{E} \asc{4}{6}{amai}{F} \asc{4}{7}{amai}{G}
    \asc{4}{8}{amai}{H} \asc{4}{9}{amai}{I} \asc{4}{10}{amai}{J} \asc{4}{11}{amai}{K}
    \asc{4}{12}{amai}{L} \asc{4}{13}{amai}{M} \asc{4}{14}{amai}{N} \asc{4}{15}{amai}{O}
    \asc{5}{0}{amai}{P} \asc{5}{1}{amai}{Q} \asc{5}{2}{amai}{R} \asc{5}{3}{amai}{S}
    \asc{5}{4}{amai}{T} \asc{5}{5}{amai}{U} \asc{5}{6}{amai}{V} \asc{5}{7}{amai}{W}
    \asc{5}{8}{amai}{X} \asc{5}{9}{amai}{Y} \asc{5}{10}{amai}{Z} \asc{5}{11}{apun}{[}
    \asc{5}{12}{apun}{\textbackslash} \asc{5}{13}{apun}{]} \asc{5}{14}{apun}{\textasciicircum} \asc{5}{15}{apun}{\_}
    \asc{6}{0}{apun}{`} \asc{6}{1}{amin}{a} \asc{6}{2}{amin}{b} \asc{6}{3}{amin}{c}
    \asc{6}{4}{amin}{d} \asc{6}{5}{amin}{e} \asc{6}{6}{amin}{f} \asc{6}{7}{amin}{g}
    \asc{6}{8}{amin}{h} \asc{6}{9}{amin}{i} \asc{6}{10}{amin}{j} \asc{6}{11}{amin}{k}
    \asc{6}{12}{amin}{l} \asc{6}{13}{amin}{m} \asc{6}{14}{amin}{n} \asc{6}{15}{amin}{o}
    \asc{7}{0}{amin}{p} \asc{7}{1}{amin}{q} \asc{7}{2}{amin}{r} \asc{7}{3}{amin}{s}
    \asc{7}{4}{amin}{t} \asc{7}{5}{amin}{u} \asc{7}{6}{amin}{v} \asc{7}{7}{amin}{w}
    \asc{7}{8}{amin}{x} \asc{7}{9}{amin}{y} \asc{7}{10}{amin}{z} \asc{7}{11}{apun}{\{}
    \asc{7}{12}{apun}{|} \asc{7}{13}{apun}{\}} \asc{7}{14}{apun}{\textasciitilde} \asc{7}{15}{actrl}{\ascnome{DEL}}
    % legenda
    \foreach \sty/\testo/\xx in {actrl/{caratteri di controllo (non stampabili)}/-0.6,
                               apun/{spazio e simboli}/4.6, adig/cifre/8.0,
                               amai/maiuscole/10.3, amin/minuscole/12.9} {
      \node[acella, \sty, minimum width=0.5cm] at (\xx, -6.85) {};
      \node[font=\Lato\scriptsize, anchor=west] at (\xx+0.35, -6.85) {\testo};
    }
  \end{tikzpicture}
  \caption{La tabella ASCII: accanto a ogni codice (da 0 a 127) il carattere
           corrispondente. I caratteri di controllo hanno un nome
           convenzionale, o la sequenza di escape con cui si scrivono in C;
           \code{\char"2423} è lo spazio.}
  \label{fig:ascii}
\end{figure}

La tabella non va imparata a memoria: in C si scrive \code{'A'} e il
compilatore lo sostituisce con 65. Conviene però capirne la struttura:

\begin{itemize}
  \item i codici da 0 a 31 e il 127 sono \textbf{caratteri di controllo}:
        non si stampano, ma comandano il dispositivo (il 10, per esempio, è
        l'a capo); il 32 è lo spazio;
  \item la scelta dei numeri è una \textbf{convenzione}: nella lettera
        \textit{A} non c'è nulla che la leghi a 65;
  \item però le cifre, le maiuscole e le minuscole occupano ciascuna dei
        codici \textbf{consecutivi} e in ordine: \code{'0'} < \code{'1'} <
        \dots{} < \code{'9'}, \code{'a'} < \code{'b'} < \dots{} <
        \code{'z'}. Confrontare i codici equivale quindi a confrontare
        nell'ordine alfabetico, e questo servirà per ordinare i testi;
  \item le maiuscole vengono prima delle minuscole, e ogni minuscola dista
        dalla sua maiuscola esattamente 32 posizioni: \code{'a'} vale 97,
        \code{'A'} vale 65;
  \item il carattere \code{'7'} non è il numero 7: vale 55. Il numero si
        ottiene con la differenza \code{'7' - '0'}, cioè $55 - 48 = 7$.
\end{itemize}

\info{E le lettere accentate?}{
  La tabella ASCII non contiene \textit{è}, \textit{à}, \textit{ù}. I sistemi
  di oggi (Linux, macOS, VS Code) usano la codifica \textbf{UTF-8}, che
  estende l'ASCII: i primi 128 codici sono gli stessi, mentre gli altri
  caratteri occupano da 2 a 4 byte. Una lettera accentata, quindi,
  \textbf{non sta in un \code{char}}: \code{char c = 'è';} produce due avvisi
  (\textit{multi-character character constant} e \textit{overflow in
  conversion}) e un valore senza senso. Dentro una stringa, invece, non c'è
  problema: \code{printf("Perché?\textbackslash n");} stampa correttamente,
  perché la stringa contiene semplicemente i due byte della \textit{é}.}

\subsection{Caratteri come interi, interi come caratteri}

Poiché un \code{char} contiene un numero, lo si può stampare in due modi:
con \code{\%c} \code{printf} disegna il carattere che ha quel codice, con
\code{\%d} scrive il numero. E siccome è un numero, lo si può usare nelle
espressioni aritmetiche: \code{c + 1} è il codice del carattere successivo.

\begin{affianco}[0.56][c]
\begin{lstlisting}[style=c]
char c = 'a';

printf("%c\n", c);
printf("%d\n", c);
printf("%c%c%c\n", c, c + 1, c + 2);
\end{lstlisting}
\risultato
\begin{console}[Terminale]
a
97
abc
\end{console}
\end{affianco}

Con la stessa idea si scorrono i caratteri con un ciclo: la variabile di
controllo parte da un codice e arriva a un altro. Nel programma seguente il
primo ciclo usa un \code{int} che va da \code{'a'} a \code{'z'}, il secondo
un \code{char} che va da 65 a 90 (i codici di \code{'A'} e \code{'Z'}), il
terzo stampa con \code{\%d} i codici delle dieci cifre.

\begin{affianco}[0.5]
\begin{lstlisting}[style=c, name=caratteri.c]
#include <stdio.h>

int main(void)
{
   int i;
   char c;

   /* stampa abc...z */
   for (i = 'a'; i <= 'z'; i++)
      printf("%c", i);
   printf("\n");

   /* stampa ABC...Z */
   for (c = 65; c <= 90; c++)
      printf("%c", c);
   printf("\n");

   /* stampa 48 49 ... 57 */
   for (c = '0'; c <= '9'; c++)
      printf("%d ", c);
   printf("\n");
   return 0;
}
\end{lstlisting}
\risultato
\begin{console}
$ ./caratteri
abcdefghijklmnopqrstuvwxyz
ABCDEFGHIJKLMNOPQRSTUVWXYZ
48 49 50 51 52 53 54 55 56 57
\end{console}
\end{affianco}

\warning{Una svista nelle slide}{
  Nelle slide il secondo ciclo è scritto \code{for (c=65; i<=90; c++)}: la
  condizione controlla \code{i} invece di \code{c}. Alla fine del primo
  ciclo \code{i} vale 123 (il primo valore per cui \code{i <= 'z'} è
  falsa), quindi \code{i <= 90} è falsa fin dall'inizio e il ciclo non stampa
  nulla. Il compilatore non può accorgersene: il codice è corretto dal punto
  di vista della sintassi, e fa semplicemente un'altra cosa.}

\section{Le sequenze di escape}

\subsection{Caratteri che non si possono scrivere}

Alcuni caratteri non si possono scrivere direttamente fra apici o fra
virgolette: o perché non sono stampabili (l'a capo, la tabulazione), o
perché hanno già un significato (le virgolette chiudono una stringa,
l'apice chiude una costante carattere). Per questi il C usa le
\term{sequenze di escape}: un backslash seguito da un carattere. Sono
\textbf{due simboli} nel codice, ma \textbf{un solo carattere} in memoria:
\code{'\textbackslash n'} è il carattere di codice 10.

\begin{center}
\renewcommand{\arraystretch}{1.12}
\begin{tabular}{@{}l>{\FiraCodeNL}c c>{\raggedright\arraybackslash}p{9.2cm}@{}}
  \toprule
  \textbf{carattere} & \normalfont\textbf{in C} & \textbf{codice} & \textbf{che cosa fa} \\
  \midrule
  a capo (\textit{newline})                   & \textbackslash n & 10 & porta il cursore all'inizio della riga successiva \\
  tabulazione (\textit{horizontal tab})       & \textbackslash t & 9  & sposta il cursore alla prossima tabulazione (ogni 8 colonne) \\
  ritorno carrello (\textit{carriage return}) & \textbackslash r & 13 & riporta il cursore all'inizio della riga \\
  tabulazione verticale (\textit{vertical tab}) & \textbackslash v & 11 & scende di una riga, restando nella stessa colonna \\
  salto pagina (\textit{formfeed})            & \textbackslash f & 12 & sulla stampante cambia pagina; nel terminale fa come \code{\textbackslash v} \\
  indietro (\textit{backspace})               & \textbackslash b & 8  & sposta il cursore indietro di una posizione \\
  campanello (\textit{alert})                 & \textbackslash a & 7  & emette un segnale acustico (o visivo) \\
  carattere nullo (\textit{null})             & \textbackslash 0 & 0  & il carattere di codice 0, che chiude le stringhe \\
  backslash                                   & \textbackslash\textbackslash & 92 & il carattere \code{\textbackslash} \\
  virgolette (\textit{double quote})          & \textbackslash" & 34 & il carattere \code{"} \\
  apice (\textit{single quote})               & \textbackslash' & 39 & il carattere \code{'} \\
  punto interrogativo                         & \textbackslash? & 63 & il carattere \code{?} (serve di rado) \\
  \bottomrule
\end{tabular}
\end{center}

\subsection{Virgolette, apici e campanello}

Per stampare delle virgolette dentro una stringa bisogna scriverle con il
backslash, altrimenti chiuderebbero la stringa. Il programma delle lezioni
lo fa in un modo un po' tortuoso, stampando con \code{\%c} il carattere
\code{'\textbackslash"'}; il modo diretto è la seconda riga:

\begin{affianco}[0.56][c]
\begin{lstlisting}[style=c]
printf("Di nome %cNemo%c\n", '\"', '\"');
printf("Di nome \"Nemo\"\n");
\end{lstlisting}
\risultato
\begin{console}
Di nome "Nemo"
Di nome "Nemo"
\end{console}
\end{affianco}

Fra apici le virgolette non danno fastidio, quindi anche \code{'"'} va bene;
viceversa, l'apice va protetto fra apici (\code{'\textbackslash'{}'}) ma non
in una stringa (\code{"l'una"}).

\code{printf("\%c", '\textbackslash a');} fa suonare il
\textbf{campanello} del terminale: un'eredità delle telescriventi, che ne
avevano uno vero. Nel terminale di VS Code, a seconda delle impostazioni, si
sente un suono, la finestra lampeggia o non succede nulla.

\subsection{Il cursore e i caratteri di controllo}

Il programma seguente stampa la stessa stringa con in mezzo, a turno,
diversi caratteri di controllo. Ecco che cosa compare sul terminale:

\begin{affianco}[0.5][c]
\begin{lstlisting}[style=c, name=escape.c]
#include <stdio.h>

int main(void)
{
  printf("abc \t abc\n");
  printf("abc \v abc\n");
  printf("abc \f abc\n");
  printf("abc \r abc\n");
  printf("abc \r \v abc\n");
  return 0;
}
\end{lstlisting}
\risultato
\begin{console}
$ ./escape
abc      abc
abc
     abc
abc
     abc
 abc
 bc
  abc
\end{console}
\end{affianco}

Per capirlo basta seguire il cursore (figura~\ref{fig:cursore}). Ogni
carattere stampabile viene scritto dove si trova il cursore, che poi avanza
di una colonna; i caratteri di controllo, invece, spostano soltanto il
cursore.

\begin{figure}[H]
  \centering
  \begin{tikzpicture}[font=\FiraCodeNL\scriptsize,
      tc/.style={mcella, minimum width=0.38cm, minimum height=0.46cm, inner xsep=0pt},
      tcn/.style={tc, fill=accent!14!white},
      tit/.style={font=\FiraCodeNL\scriptsize, anchor=west, inner sep=1pt},
      esc/.style={font=\FiraCodeNL\scriptsize\bfseries, text=accent!75!black, inner sep=1.5pt},
      lab/.style={font=\Lato\scriptsize, text=neutral!40!black, anchor=east}]
    \noindici
    % \t
    \node[tit] at (0,1.35) {printf("abc \textbackslash t abc\textbackslash n");};
    \celle[tc]{ta}{(0,0)}{a,b,c,\char"2423,{},{},{},{},\char"2423,a,b,c}
    \draw[wtarrowacc] ([yshift=1pt]ta-4.north west) to[bend left=55]
      node[esc, above] {\textbackslash t} ([yshift=1pt]ta-8.north west);
    \node[mindice, anchor=north] at ([yshift=-1pt]ta-0.south) {0};
    \node[mindice, anchor=north] at ([yshift=-1pt]ta-4.south) {4};
    \node[mindice, anchor=north] at ([yshift=-1pt]ta-8.south) {8};
    \node[mnota, anchor=north west, text width=4.4cm] at (-0.1,-0.6)
      {salta alla prossima colonna multipla di 8: qui dalla 4 alla 8};
    % \v
    \node[tit] at (6.0,1.35) {printf("abc \textbackslash v abc\textbackslash n");};
    \celle[tc]{tb}{(6.0,0)}{a,b,c,\char"2423,{},{},{},{}}
    \celle[tc]{tv}{(6.0,-0.85)}{{},{},{},{},\char"2423,a,b,c}
    \draw[wtarrowacc] (tb-4.center) -- node[esc, right=3pt, pos=0.45] {\textbackslash v} (tv-4.north);
    \node[mnota, anchor=north west, text width=4.4cm] at (5.9,-1.25)
      {scende di una riga e resta nella stessa colonna};
    % \r
    \node[tit] at (11.6,1.35) {printf("abc \textbackslash r abc\textbackslash n");};
    \celle[tc]{tr}{(12.4,0)}{a,b,c,\char"2423}
    \celle[tcn]{ts}{(12.4,-0.85)}{\char"2423,a,b,c}
    \node[lab] at (12.3,0) {prima};
    \node[lab] at (12.3,-0.85) {dopo};
    \draw[wtarrowacc] ([yshift=1pt]tr-3.north east) to[bend right=55]
      node[esc, above] {\textbackslash r} ([yshift=1pt]tr-0.north west);
    \node[mnota, anchor=north west, text width=4.4cm] at (11.5,-1.25)
      {torna all'inizio della riga: ciò che segue sovrascrive i caratteri già scritti};
  \end{tikzpicture}
  \caption{Il percorso del cursore con \code{\textbackslash t},
           \code{\textbackslash v} e \code{\textbackslash r}
           (\code{\char"2423} è uno spazio).}
  \label{fig:cursore}
\end{figure}

\begin{itemize}
  \item \code{\textbackslash t} sposta il cursore alla prossima
        \textbf{tabulazione}; le tabulazioni sono ogni 8 colonne (0, 8,
        16\dots), e per questo servono ad allineare i dati in colonna.
  \item \code{\textbackslash v} e \code{\textbackslash f} fanno scendere il
        cursore di una riga, ma senza riportarlo a sinistra.
  \item \code{\textbackslash r} riporta il cursore all'inizio della riga
        \textbf{senza} scendere: ciò che si stampa dopo sovrascrive. Così
        \code{"abc "} diventa \code{" abc"}.
  \item L'ultima riga combina le due cose: \code{\textbackslash r} torna
        all'inizio, lo spazio sovrascrive la \code{a} (si legge \code{" bc"}),
        \code{\textbackslash v} scende di una riga restando nella colonna 1,
        e lì viene scritto \code{" abc"}.
\end{itemize}

E \code{\textbackslash n}? Nel terminale fa due cose insieme: scende di una
riga \textbf{e} torna all'inizio. Sulle macchine da scrivere erano due
movimenti distinti, riportare il carrello all'inizio (\textit{carriage
return}) e far scorrere il foglio di una riga (\textit{line feed}): per
questo i file di testo di Windows terminano ogni riga con due caratteri,
\code{\textbackslash r\textbackslash n}, mentre Linux e macOS usano solo
\code{\textbackslash n}.

\subsection{Un'animazione con \code{\textbackslash r}}

Il ritorno carrello permette di ridisegnare più volte lo stesso punto dello
schermo. Il programma seguente stampa a ripetizione i caratteri
\code{|}, \code{/}, \code{-} e \code{\textbackslash}, ciascuno seguito da
\code{\textbackslash r}, con una pausa di un secondo fra l'uno e l'altro:
ogni simbolo prende il posto del precedente e si vede una stanghetta che
ruota.

\begin{affianco}[0.5][c]
\begin{lstlisting}[style=c, name=spinner.c]
#include <stdio.h>
#include <unistd.h>

#define PAUSA 1

void pausa(int l)
{
    sleep(l);
}

int main()
{
    setbuf(stdout, NULL);
    while (1) {
        printf("|\r");  pausa(PAUSA);
        printf("/\r");  pausa(PAUSA);
        printf("-\r");  pausa(PAUSA);
        printf("\\\r"); pausa(PAUSA);
    }
}
\end{lstlisting}
\risultato
\begin{itemize}
  \item \code{"\textbackslash\textbackslash\textbackslash r"} contiene due
        caratteri: il backslash (scritto \code{\textbackslash\textbackslash})
        e il ritorno carrello.
  \item \code{while (1)} è un ciclo infinito, perché la condizione è sempre
        vera: il programma si ferma con \tasti{Ctrl,C}.
  \item \code{sleep(l)}, dichiarata in \code{unistd.h}, sospende il
        programma per \code{l} secondi. È disponibile su Linux, su macOS e
        in WSL, non su Windows senza WSL.
\end{itemize}
\end{affianco}

Nella cartella degli esempi c'è anche una versione senza \code{sleep}, in
cui \code{pausa} perde tempo con un ciclo vuoto:
\code{for (i = 0; i < l; i++);} (il \code{;} da solo è l'istruzione vuota).
Funziona, ma la durata della pausa dipende dalla velocità del computer; e
compilando con l'ottimizzazione (\code{gcc -O2}) il compilatore si accorge
che il ciclo non fa nulla e lo elimina: in due secondi la versione ottimizzata
ha stampato oltre 2,8 milioni di caratteri, quella normale 70.

\subsection{Il buffer di \code{printf}}

Resta da capire \code{setbuf(stdout, NULL)}. Per efficienza, \code{printf}
non manda ogni carattere allo schermo appena lo riceve: lo accumula in una
zona di memoria, il \term{buffer}, e la svuota tutta insieme quando incontra
un \code{\textbackslash n}, quando è piena o quando il programma termina
(figura~\ref{fig:buffer}). Il nostro programma non stampa mai un
\code{\textbackslash n} e non termina mai: senza \code{setbuf} il buffer non
verrebbe mai svuotato e sullo schermo \textbf{non comparirebbe nulla}.
L'abbiamo verificato: senza quella riga, in tre secondi il programma non
mostra niente. \code{setbuf(stdout, NULL)} elimina il buffer dell'output
standard (\code{stdout}), e ogni carattere arriva subito allo schermo.

\begin{figure}[H]
  \centering
  \begin{tikzpicture}[font=\Lato\scriptsize]
    \noindici
    \node[font=\Lato\scriptsize\bfseries, anchor=west] at (-0.9,0.85) {senza \code{setbuf}};
    \node[wtnode, font=\Lato\scriptsize] (p1) at (0,0) {programma};
    \celle[mcella, minimum width=0.55cm]{bf}{(1.7,0)}{|,\textbackslash r,/,\textbackslash r,-,\textbackslash r,{}}
    \node[mnota, anchor=south] at (bf-box.north) {buffer di \code{stdout}: i caratteri si accumulano};
    \node[wtnodeplain] (s1) at (9.2,0) {schermo};
    \draw[wtarrow] (p1) -- (bf-box.west);
    \draw[wtarrow, dashed, draw=merr] (bf-box.east) -- node[above, mnota, text=merr] {solo con \code{\textbackslash n}} (s1);
    \node[font=\Lato\scriptsize\bfseries, anchor=west] at (-0.9,-0.95) {con \code{setbuf(stdout, NULL)}};
    \node[wtnode, font=\Lato\scriptsize] (p2) at (0,-1.8) {programma};
    \node[wtnodeplain] (s2) at (9.2,-1.8) {schermo};
    \draw[wtarrow, draw=mok] (p2) -- node[above, mnota] {ogni carattere arriva subito} (s2);
  \end{tikzpicture}
  \caption{Con il buffer l'output arriva allo schermo a righe intere; senza,
           un carattere alla volta.}
  \label{fig:buffer}
\end{figure}

\section{Attenzione a \code{++} dentro \code{printf}}

Le slide mettono in guardia: con gli operatori \code{++} e \code{--} a volte
i risultati \textbf{non sono prevedibili}. Il programma seguente stampa più
volte \code{c} e \code{a} incrementandoli dentro la stessa chiamata di
\code{printf}. Compilato con \code{gcc} dà i risultati delle slide; compilato
con \code{clang} dà risultati del tutto diversi:

\begin{affianco}[0.5][t]
\begin{lstlisting}[style=c, name=printf++.c]
#include <stdio.h>

int main()
{
  char c;  int a;

  c = 'a'; a=10;
  printf("%c%c%c\n", c, c, c++);
  printf ("%d %d %d \n", a , a , a++);

  c = 'a'; a=10;
  printf("%c%c%c\n", c, c, ++c);
  printf("%d %d %d \n", a , a , ++a);

  c = 'a'; a=10;
  printf("%c%c%c\n", c, ++c, ++c);
  printf("%d %d %d \n", a , ++a , ++a);

  c = 'a'; a=10;
  printf("%c%c%c\n", c, c++, c++);
  printf("%d %d %d \n", a , a++ , a++);

  c = 'a'; a=10;
  printf("%c%c%c \n", ++c , c++ , c);
  printf("%d %d %d \n", ++a , a++ , a);
}
\end{lstlisting}
\risultato
\begin{minipage}[t]{0.48\linewidth}
\begin{console}[gcc]
bba
11 11 10
bbb
11 11 11
ccb
12 12 12
cba
12 11 10
caa
12 10 12
\end{console}
\end{minipage}\hfill
\begin{minipage}[t]{0.48\linewidth}
\begin{console}[clang]
aaa
10 10 10
aab
10 10 11
abc
10 11 12
aab
10 10 11
bbc
11 11 12
\end{console}
\end{minipage}
\end{affianco}

Nessuno dei due compilatori sbaglia. Il C non stabilisce in che ordine
vengano calcolati gli argomenti di una funzione (qui \code{gcc} li ha
calcolati da destra a sinistra, \code{clang} da sinistra a destra); e
modificare una variabile con \code{++} e usarla di nuovo nella stessa
chiamata dà un comportamento \textbf{indefinito}, come \code{a * a++} nel
capitolo~\ref{chap:sintassi}. Con \code{-Wall}, \code{gcc} lo segnala per
ogni riga:

\begin{console}
$ gcc -Wall -o printf++ printf++.c
printf++.c: In function (*@‘@*)main(*@’@*):
printf++.c:8:29: warning: operation on (*@‘@*)c(*@’@*) may be undefined [-Wsequence-point]
    8 |   printf("%c%c%c\n", c, c, c++);
      |                            ~^~
\end{console}

La regola pratica: una variabile incrementata con \code{++} o decrementata
con \code{--} \textbf{non va usata altrove} nella stessa espressione o nella
stessa chiamata. Si incrementa prima, in un'istruzione a parte.

\section{Come è fatto un \code{char}}

Vediamo ora che cosa c'è davvero nel byte di un \code{char}, e perché
esistono tre tipi diversi per i caratteri.

\subsection{I bit di un carattere}

La notazione binaria (capitolo~\ref{chap:calcolatore}) è
\textbf{posizionale} come quella decimale. In base 10 il numero
$d_n d_{n-1} \dots d_1 d_0$ vale
$d_n \cdot 10^n + \dots + d_1 \cdot 10^1 + d_0 \cdot 10^0$ (per esempio
$20753 = 2 \cdot 10^4 + 0 \cdot 10^3 + 7 \cdot 10^2 + 5 \cdot 10^1 + 3 \cdot
10^0$: il 3 è la cifra meno significativa, il 2 la più significativa). In
base 2 ogni cifra $b_i$ è un bit, 0 o 1, e il numero vale
$b_n \cdot 2^n + \dots + b_1 \cdot 2^1 + b_0 \cdot 2^0$. Dopo la
dichiarazione \code{char c = 'a';} il byte di \code{c} contiene
\code{01100001} (figura~\ref{fig:bit-a}).

\begin{figure}[H]
  \centering
  \begin{tikzpicture}[font=\FiraCodeNL\small]
    \foreach \p [count=\i from 0] in {128,64,32,16,8,4,2,1} {
      \node[mcella, minimum width=1.1cm, minimum height=0.65cm] (b\i) at (\i*1.1,0) {};
      \node[font=\FiraCodeNL\scriptsize, text=neutral!40!black] at (\i*1.1,0.58) {\p};
      \pgfmathtruncatemacro{\e}{7-\i}
      \node[font=\Lato\tiny, text=neutral!50!black] at (\i*1.1,0.95) {$2^{\e}$};
    }
    \foreach \val [count=\i from 0] in {0,1,1,0,0,0,0,1} {
      \ifnum\val=1
        \node[font=\FiraCodeNL\normalsize\bfseries, text=accent!80!black] at (b\i) {1};
      \else
        \node[font=\FiraCodeNL\normalsize, text=neutral!55!black] at (b\i) {0};
      \fi
    }
    \node[anchor=west, font=\Lato\small] at (8.4,0)
      {$= 64 + 32 + 1 = \mathbf{97}$, il codice di \code{'a'}};
  \end{tikzpicture}
  \caption{Il byte del carattere \code{'a'}: ogni bit a 1 aggiunge il peso
           della sua posizione.}
  \label{fig:bit-a}
\end{figure}

\subsection{Tre tipi per i caratteri}

Il C ha tre tipi da un byte: \code{char}, \code{signed char} e
\code{unsigned char}. Ognuno ha 256 valori distinti, ma non gli stessi:

\begin{itemize}
  \item \code{signed char} va da $-128$ a $127$;
  \item \code{unsigned char} va da $0$ a $255$;
  \item \code{char} è uno dei due, a seconda dell'\textbf{implementazione}
        (cioè del compilatore e del processore).
\end{itemize}

Per i caratteri ASCII (da 0 a 127) la differenza non conta. Conta quando il
byte si usa come numero: per capirla bisogna vedere come si rappresentano i
numeri negativi.

\subsection{I numeri negativi}

L'idea più semplice è il \textbf{bit di segno}. Limitiamoci a 4 bit: il
bit più significativo indica il segno e gli altri tre il valore, quindi
$5 = $~\code{0101} e $-5 = $~\code{1101}. Lo zero però ha due
rappresentazioni, \code{0000} e \code{1000}, e così un valore va sprecato.

Per questo i computer usano il \term{complemento a due}: i numeri positivi si
rappresentano normalmente, e per un negativo si \textbf{invertono} i bit del
corrispondente positivo e si \textbf{somma 1}. Così \code{1000} non è più
uno zero, ma $-8$ (figura~\ref{fig:complemento}).

\begin{figure}[H]
  \centering
  \begin{tikzpicture}[font=\FiraCodeNL\scriptsize,
      vn/.style={mcella, minimum width=0.95cm, font=\FiraCodeNL\scriptsize\bfseries, fill=accent!10!white},
      vp/.style={mcella, minimum width=0.95cm, font=\FiraCodeNL\scriptsize\bfseries},
      pn/.style={mcella, minimum width=0.95cm, fill=accent!10!white},
      pp/.style={mcella, minimum width=0.95cm},
      lab/.style={font=\Lato\scriptsize, anchor=east, text=neutral!40!black}]
    \noindici
    \celle[vn]{vn}{(0,0)}{$-8$,$-7$,$-6$,$-5$,$-4$,$-3$,$-2$,$-1$}
    \celle[vp]{vp}{(vn-7.east)}{0,1,2,3,4,5,6,7}
    \celle[pn]{pn}{(0,-0.5)}{1000,1001,1010,1011,1100,1101,1110,1111}
    \celle[pp]{pp}{(pn-7.east)}{0000,0001,0010,0011,0100,0101,0110,0111}
    \node[lab] at (-0.1,0) {valore};
    \node[lab] at (-0.1,-0.5) {bit};
    \draw[wtbrace] ([yshift=-3pt]pn-0.south west) -- ([yshift=-3pt]pn-7.south east);
    \draw[wtbrace] ([yshift=-3pt]pp-0.south west) -- ([yshift=-3pt]pp-7.south east);
    \node[mnota, anchor=north] at ([yshift=-0.32cm]$(pn-0.south)!0.5!(pn-7.south)$) {il primo bit è 1: numeri negativi};
    \node[mnota, anchor=north] at ([yshift=-0.32cm]$(pp-0.south)!0.5!(pp-7.south)$) {il primo bit è 0: zero e numeri positivi};
  \end{tikzpicture}
  \caption{I 16 valori di 4 bit in complemento a due: da $-8$ a $+7$, con
           un solo zero.}
  \label{fig:complemento}
\end{figure}

\begin{affianco}[0.36][c]
\begin{center}
\renewcommand{\arraystretch}{1.15}
\begin{tabular}{@{}l>{\FiraCodeNL}r@{}}
  \toprule
  $5$ & 0101 \\
  bit invertiti & 1010 \\
  sommo 1 & + 0001 \\
  \midrule
  $-5$ & 1011 \\
  \bottomrule
\end{tabular}
\end{center}
\risultato
Per esempio, $-5$ si ottiene da \code{0101}: invertendo i bit si ha
\code{1010}, e sommando 1 si arriva a \code{1011}. Perché funziona? Un numero
più il suo inverso dà sempre \code{1111} (in ogni colonna c'è un 1 e uno 0,
senza riporti), cioè $-1$. Quindi l'inverso di $x$ vale $-1 - x$, e
aggiungendo 1 si ottiene proprio $-x$.
\end{affianco}

\subsection{Lo stesso byte, due letture}

Su 8 bit funziona allo stesso modo. I byte da \code{00000000} a
\code{01111111} hanno lo stesso valore, da 0 a 127, per \code{signed char} e
per \code{unsigned char}. Quelli che cominciano con un 1 no: per
\code{unsigned char} vanno da 128 a 255, per \code{signed char} da $-128$ a
$-1$ (figura~\ref{fig:due-letture}). Il byte \code{11111111}, per esempio,
vale 255 oppure $-1$.

\begin{figure}[H]
  \centering
  \begin{tikzpicture}[font=\FiraCodeNL\scriptsize,
      lab/.style={font=\Lato\scriptsize\bfseries, anchor=east}]
    \fill[green-gradient-1!40!white] (0,0) rectangle (6,0.55);
    \fill[accent!12!white] (6,0) rectangle (12,0.55);
    \draw[primary!65!black, line width=0.7pt] (0,0) rectangle (12,0.55) (6,0) -- (6,0.55);
    \node at (3,0.275) {00000000 \ldots\ 01111111};
    \node at (9,0.275) {10000000 \ldots\ 11111111};
    \node[lab] at (-0.35,0.275) {byte};
    \node[lab] at (-0.35,1.05) {unsigned char};
    \node[lab] at (-0.35,-0.5) {signed char};
    \foreach \xx/\val in {0.2/0, 5.7/127, 6.3/128, 11.8/255}
      \node at (\xx,1.05) {\val};
    \foreach \xx/\val in {0.2/0, 5.7/127, 6.3/$-128$, 11.8/$-1$}
      \node at (\xx,-0.5) {\val};
    \draw[wtarrow] (0.6,1.05) -- (5.25,1.05);
    \draw[wtarrow] (6.75,1.05) -- (11.4,1.05);
    \draw[wtarrow] (0.6,-0.5) -- (5.25,-0.5);
    \draw[wtarrow, draw=accent!70!black] (6.85,-0.5) -- (11.35,-0.5);
    \node[mnota, align=center, anchor=north] at (3,-0.75) {le due letture coincidono};
    \node[mnota, align=center, anchor=north] at (9,-0.75) {le due letture sono diverse};
  \end{tikzpicture}
  \caption{I 256 byte letti come \code{unsigned char} (sopra) e come
           \code{signed char} (sotto).}
  \label{fig:due-letture}
\end{figure}

\subsection{Il \code{char} del mio computer ha il segno?}

Come si scopre se \code{char} è \code{signed} o \code{unsigned}? Si assegna
$-1$ ai tre tipi e si stampano i valori con \code{\%d}:

\begin{affianco}[0.5][c]
\begin{lstlisting}[style=c, name=char.c]
#include <stdio.h>

int main()
{
    char c = -1;
    unsigned char u = -1;
    signed char s = -1;

    printf("c = %d  u = %d   s = %d\n", c, u, s);

    return 0;
}
\end{lstlisting}
\risultato
\begin{console}
$ gcc -Wall -o char char.c
$ ./char
c = -1  u = 255   s = -1
$ gcc -Wall -funsigned-char -o char char.c
$ ./char
c = 255  u = 255   s = -1
\end{console}
\end{affianco}

In tutte e tre le variabili il byte è \code{11111111}: $-1$ assegnato a un
\code{unsigned char}, che non ha negativi, diventa $-1 + 256 = 255$. La
differenza sta nella lettura: per stampare, il valore viene convertito in
\code{int} e conserva il suo significato, quindi $-1$ per il tipo con segno e
255 per quello senza. Sui PC (processori Intel e AMD) \code{char} ha il
segno, e infatti \code{c} vale $-1$. Su altri sistemi, per esempio Linux sui
processori ARM come quelli del Raspberry Pi, \code{char} è senza segno; lo si
può simulare con l'opzione \code{-funsigned-char} di \code{gcc}, ed ecco
\code{c = 255}.

La conclusione pratica: per i caratteri si usa \code{char}; quando un byte
serve come \textbf{numero} piccolo e il segno conta, si scrive esplicitamente
\code{signed char} o \code{unsigned char}.

\section{Il tipo \code{int}}

\code{int} è il tipo principale del C, quello che si usa per qualunque
numero intero quando non ci sono motivi per sceglierne un altro.

\subsection{Quanti valori}

Un \code{int} occupa 2 byte (16 bit) sui sistemi vecchi e 4 byte (32 bit)
su quasi tutti i sistemi moderni: il numero di valori distinti, quindi,
dipende dal sistema. In complemento a due, con $n$ bit si rappresentano gli
interi da $-2^{n-1}$ a $2^{n-1}-1$:

\begin{itemize}
  \item con 16 bit: da $-2^{15} = -32\,768$ a $2^{15} - 1 = 32\,767$;
  \item con 32 bit: da $-2^{31} = -2\,147\,483\,648$ a
        $2^{31} - 1 = 2\,147\,483\,647$, circa più o meno 2 miliardi.
\end{itemize}

Il massimo è $2^{n-1}-1$ e non $2^{n-1}$ perché fra i valori con il primo
bit a 0 c'è anche lo zero. Le slide chiamano questi estremi $N_{min\_int}$ e
$N_{max\_int}$: ogni variabile \code{int i} soddisfa
$N_{min\_int} \le i \le N_{max\_int}$. In C sono disponibili come costanti
simboliche, \code{INT\_MIN} e \code{INT\_MAX}, definite nell'header
\code{limits.h}.

\subsection{L'overflow}

Che cosa succede se un calcolo supera \code{INT\_MAX}? Il programma seguente
somma due miliardi a due miliardi:

\begin{affianco}[0.5][c]
\begin{lstlisting}[style=c, name=overflow.c]
#include <stdio.h>

#define BIG 2000000000   /* 2 miliardi */

int main(void)
{
   int a, b = BIG, c = BIG;

   a = b + c;
   printf("a = %d\n", a);
   return 0;
}
\end{lstlisting}
\risultato
\begin{console}
$ gcc -Wall -o overflow overflow.c
$ ./overflow
a = -294967296
\end{console}
\end{affianco}

Il codice è sintatticamente corretto e il compilatore non segnala nulla; ma
in esecuzione $b + c$ vale 4 miliardi, più di \code{INT\_MAX}: è un
\term{integer overflow} (traboccamento). Il programma prosegue come se
niente fosse, con un risultato sbagliato. La figura~\ref{fig:overflow} mostra
da dove viene proprio quel numero: superato il massimo si ricomincia dal
minimo, come il contachilometri di un'auto che dopo 99999 torna a 00000 (la
ruota del capitolo~\ref{chap:calcolatore}). Contando, si arriva a
$4\,000\,000\,000 - 2^{32} = -294\,967\,296$.

\begin{figure}[H]
  \centering
  \begin{tikzpicture}[font=\Lato\scriptsize]
    \draw[line width=1.2pt, neutral!45!black] (0,0) -- (12,0);
    \foreach \xx/\lab in {0/{\code{INT\_MIN}\\$-2\,147\,483\,648$}, 6/{0},
                          12/{\code{INT\_MAX}\\$2\,147\,483\,647$}} {
      \draw[line width=0.8pt, neutral!45!black] (\xx,-0.12) -- (\xx,0.12);
      \node[anchor=north, align=center] at (\xx,-0.2) {\lab};
    }
    \fill[primary!70!black] (11.59,0) circle (2.3pt);
    \node[anchor=north, align=center, text=primary!60!black] at (10.2,-0.2) {\code{b}\\$2\,000\,000\,000$};
    \draw[primary!50!black, line width=0.5pt] (11.59,-0.06) -- (10.6,-0.25);
    \draw[wtarrowacc] (11.59,0.25) -- node[above, text=accent!70!black] {$+$ \code{c}} (12.7,0.25);
    \draw[accent!70!black, dashed, line width=0.9pt] (12.7,0.25) .. controls (14.2,2.3) and (-2.2,2.3) .. (-0.7,0.25);
    \node[fill=white, inner sep=2pt, text=accent!70!black] at (6,1.75) {superato \code{INT\_MAX}, si ricomincia da \code{INT\_MIN}};
    \draw[wtarrowacc] (-0.7,0.25) -- (5.18,0.25);
    \fill[accent!80!black] (5.18,0) circle (2.3pt);
    \node[anchor=north, align=center, text=accent!70!black] at (4.6,-0.2) {\code{a}\\$-294\,967\,296$};
  \end{tikzpicture}
  \caption{\code{b + c} su 32 bit: la somma esce a destra e rientra da
           sinistra.}
  \label{fig:overflow}
\end{figure}

Il risultato si spiega così: la somma viene fatta su 32 bit, e i bit che non
ci stanno si perdono. Per i tipi \code{unsigned} questo “giro” è garantito
dallo standard (i calcoli sono fatti modulo $2^n$); per i tipi con segno lo
standard parla invece di comportamento \textbf{indefinito}, e il valore
ottenuto qui è solo ciò che succede in pratica sui nostri computer. In ogni
caso tocca al programmatore assicurarsi che i valori restino
nell'intervallo, per esempio usando un tipo più grande.

È quello che fa la versione di \code{overflow.c} nella cartella degli
esempi: usa \code{unsigned long} (8 byte) e numeri ancora più grandi, 40
miliardi, e la somma, 80 miliardi, viene giusta.

\begin{affianco}[0.5][c]
\begin{lstlisting}[style=c, name=overflow.c]
#include <stdio.h>

#define BIG	40000000000

int main()
{
	unsigned long int a, b, c;

	b = c = BIG;

	a = b + c;

	printf("a = %ld\n", a);

	return 0;
}
\end{lstlisting}
\risultato
\begin{console}
$ ./overflow
a = 80000000000
\end{console}
(Per un \code{unsigned long} la conversione esatta di \code{printf} sarebbe
\code{\%lu}; qui \code{\%ld} funziona perché il valore sta anche in un
\code{long}.)
\end{affianco}

\section{\code{short}, \code{long} e \code{unsigned}}

\subsection{Più piccoli, più grandi, senza segno}

Oltre a \code{char} e \code{int}, gli altri tipi interi servono per usi più
specifici:

\begin{itemize}
  \item \code{short} occupa di solito meno memoria di un \code{int} (2
        byte): si usa quando lo spazio conta, per esempio in tabelle con
        milioni di numeri piccoli;
  \item \code{long} si usa per valori grandi (8 byte su Linux e macOS a 64
        bit);
  \item \code{unsigned} occupa quanto un \code{int}, ma non ha negativi: con
        4 byte va da 0 a $2^{32} - 1 = 4\,294\,967\,295$.
\end{itemize}

I limiti esatti del proprio sistema si leggono in \code{limits.h}. Questo
programma li stampa (l'output è quello di Linux su un PC a 64 bit):

\begin{lstlisting}[style=c, name=limiti.c]
#include <stdio.h>
#include <limits.h>

int main(void)
{
   printf("signed char:   %d .. %d\n", SCHAR_MIN, SCHAR_MAX);
   printf("unsigned char: %d .. %d\n", 0, UCHAR_MAX);
   printf("char:          %d .. %d\n", CHAR_MIN, CHAR_MAX);
   printf("short:         %d .. %d\n", SHRT_MIN, SHRT_MAX);
   printf("int:           %d .. %d\n", INT_MIN, INT_MAX);
   printf("unsigned:      %u .. %u\n", 0u, UINT_MAX);
   printf("long:          %ld .. %ld\n", LONG_MIN, LONG_MAX);
   printf("unsigned long: %lu .. %lu\n", 0ul, ULONG_MAX);
   return 0;
}
\end{lstlisting}

\begin{console}
$ ./limiti
signed char:   -128 .. 127
unsigned char: 0 .. 255
char:          -128 .. 127
short:         -32768 .. 32767
int:           -2147483648 .. 2147483647
unsigned:      0 .. 4294967295
long:          -9223372036854775808 .. 9223372036854775807
unsigned long: 0 .. 18446744073709551615
\end{console}

\info{Che cosa garantisce lo standard}{
  Lo standard del C non fissa le dimensioni dei tipi, ma solo dei
  \textbf{minimi}: \code{char} ha almeno 8 bit, \code{short} e \code{int}
  almeno 16, \code{long} almeno 32; inoltre un \code{short} non è più grande
  di un \code{int}, che non è più grande di un \code{long}. Le dimensioni
  reali le sceglie l'implementazione: su Windows, per esempio, \code{long}
  ha 4 byte come \code{int}, mentre su Linux e macOS a 64 bit ne ha 8. Chi
  ha bisogno di dimensioni esatte trova nell'header \code{stdint.h} (C99)
  tipi come \code{int32\_t} e \code{int64\_t}.}

\subsection{Il tipo delle costanti intere}

Anche le costanti hanno un tipo. A una costante intera scritta in decimale
viene dato il primo tipo in grado di contenerla fra \code{int} e \code{long}
(e, se non basta, \code{long long}, un tipo introdotto dal C99). Il tipo si
può anche imporre con un \textbf{suffisso}: \code{U} per \code{unsigned},
\code{L} per \code{long}, \code{UL} per \code{unsigned long} (vanno bene
anche le minuscole, ma la \code{l} si confonde facilmente con un 1). Per i
reali, come vedremo, i suffissi sono \code{F} e \code{L}. Ecco i tipi
assegnati dal compilatore ad alcune costanti:

\begin{center}
\renewcommand{\arraystretch}{1.12}
\begin{tabular}{@{}>{\FiraCodeNL}l l l@{}}
  \toprule
  \normalfont\textbf{costante} & \textbf{tipo} & \textbf{perché} \\
  \midrule
  37         & \code{int}           & sta in un \code{int} \\
  37U        & \code{unsigned}      & suffisso \code{U} \\
  37L        & \code{long}          & suffisso \code{L} \\
  37UL       & \code{unsigned long} & suffissi \code{U} e \code{L} \\
  2147483647 & \code{int}           & è \code{INT\_MAX}: sta ancora in un \code{int} \\
  2147483648 & \code{long}          & non sta più in un \code{int} \\
  3.7        & \code{double}        & le costanti reali sono \code{double} \\
  3.7F       & \code{float}         & suffisso \code{F} \\
  3.7L       & \code{long double}   & suffisso \code{L} \\
  'a'        & \code{int}           & le costanti carattere sono \code{int} \\
  \bottomrule
\end{tabular}
\end{center}

(Le slide indicano come candidati \code{int}, \code{long} e \code{unsigned}:
era all'incirca la regola del C più vecchio, il C89. Nel C99 una costante
decimale senza suffisso non diventa mai \code{unsigned}; possono diventarlo
quelle scritte in esadecimale o in ottale.)

\section{I tipi reali}

\subsection{\code{float}, \code{double} e \code{long double}}

L'ANSI C fornisce tre tipi reali: \code{float}, \code{double} e
\code{long double}. Le costanti reali, come \code{3.7}, sono di tipo
\code{double}; il tipo si può imporre con un suffisso: \code{3.7F} è un
\code{float}, \code{3.7L} un \code{long double}.

Di solito il compilatore dà più memoria ai \code{double} che ai
\code{float}, ma non è obbligatorio; su molte macchine un \code{float}
occupa 4 byte e un \code{double} 8. Allo stesso modo un \code{long double}
può avere più memoria di un \code{double} (sui PC ne occupa 16), ma non è
obbligatorio.

I valori che un tipo reale può contenere si descrivono con due grandezze:

\begin{itemize}
  \item la \textbf{precisione}: quante cifre significative, in notazione
        decimale, si possono rappresentare;
  \item l'\textbf{intervallo} (\textit{range}): il valore più grande e il più
        piccolo che si possono rappresentare.
\end{itemize}

\subsection{La virgola mobile}

Su molte macchine un \code{float} ha una precisione di 6 cifre significative
e un intervallo che va da circa $10^{-38}$ a $10^{38}$. Si può pensarlo come
un numero della forma
\[
  0.d_1 d_2 d_3 d_4 d_5 d_6 \times 10^{n}, \qquad -38 \le n \le 38,
\]
dove ogni $d_i$ è una cifra decimale (figura~\ref{fig:virgola-mobile}). Le
cifre stabiliscono la precisione, l'esponente $n$ dice dove va la virgola e
i suoi limiti stabiliscono l'intervallo. La virgola “galleggia” a seconda
dell'esponente: con le stesse sei cifre si rappresenta $123.451$ (esponente
3), $0.000123451$ (esponente $-3$) o $123\,451\,000$ (esponente 9). Da qui il
nome \textit{floating point}, virgola mobile.

\begin{figure}[H]
  \centering
  \begin{tikzpicture}[font=\FiraCodeNL\small]
    \noindici
    \node[anchor=east, font=\FiraCodeNL\small] at (-0.05,0) {123.451\Lato\ \ =\ \ \FiraCodeNL 0.};
    \celle[mcella, minimum width=0.6cm, minimum height=0.6cm, font=\FiraCodeNL\small\bfseries]{m}{(0,0)}{1,2,3,4,5,1}
    \node[anchor=west, font=\Lato\small] (x) at ([xshift=2pt]m-5.east) {$\times\ 10$};
    \node[mcella, anchor=south west, minimum width=0.45cm, minimum height=0.42cm,
          font=\FiraCodeNL\scriptsize\bfseries, fill=accent!12!white, draw=accent!65!black] (e)
      at ([xshift=-1pt, yshift=-3pt]x.north east) {3};
    \draw[wtbrace] ([yshift=-3pt]m-0.south west) -- ([yshift=-3pt]m-5.south east);
    \node[mnota, anchor=north, align=center] at ([yshift=-0.35cm]$(m-0.south)!0.5!(m-5.south)$)
      {le cifre significative:\\quante sono decide la \textbf{precisione}};
    \node[mnota, anchor=west, align=left] at ([xshift=0.35cm]e.east)
      {l'esponente dice dove va la virgola;\\i suoi limiti decidono l'\textbf{intervallo}};
  \end{tikzpicture}
  \caption{Un \code{float} come sei cifre più un esponente.}
  \label{fig:virgola-mobile}
\end{figure}

In realtà la macchina usa la base 2, non la base 10: le cifre sono binarie
(24 per un \code{float}), e “6 cifre decimali” è il modo di dire quanta
precisione garantiscono. I limiti esatti sono nell'header \code{float.h}:

\begin{lstlisting}[style=c, name=reali.c]
#include <stdio.h>
#include <float.h>

int main(void)
{
   printf("float:  min %e  max %e  cifre %d\n",
          FLT_MIN, FLT_MAX, FLT_DIG);
   printf("double: min %e  max %e  cifre %d\n",
          DBL_MIN, DBL_MAX, DBL_DIG);
   return 0;
}
\end{lstlisting}

\begin{console}
$ ./reali
float:  min 1.175494e-38  max 3.402823e+38  cifre 6
double: min 2.225074e-308  max 1.797693e+308  cifre 15
\end{console}

Un \code{double}, su molte macchine, ha quindi circa 15 cifre significative
e un intervallo da circa $10^{-308}$ a $10^{308}$: si può pensare come
$0.d_1 d_2 \dots d_{15} \times 10^{n}$ con $-308 \le n \le 308$.

\subsection{Ai confini del \code{float}}

Il programma delle lezioni assegna a un \code{float} valori sempre più
piccoli, poi due valori molto grandi e infine un numero con molte cifre:

\begin{affianco}[0.56][t]
\codicepiccolo
\begin{lstlisting}[style=c, name=float.c]
#include <stdio.h>

int main()
{
   float x;

   x = 1e-37; printf("x = 1e-37    x = %e\n", x);
   x = 1e-38; printf("x = 1e-38    x = %e\n", x);
   x = 1e-39; printf("x = 1e-39    x = %e\n", x);
   x = 1e-40; printf("x = 1e-40    x = %e\n", x);
   x = 1e-41; printf("x = 1e-41    x = %e\n", x);
   x = 1e-42; printf("x = 1e-42    x = %e\n", x);
   x = 1e-43; printf("x = 1e-43    x = %e\n", x);
   x = 1e-44; printf("x = 1e-44    x = %e\n", x);
   x = 1e-45; printf("x = 1e-45    x = %e\n", x);
   x = 1e-46; printf("x = 1e-46    x = %e\n", x);

   x = 1e38; printf("x = 1e38    x = %e\n", x);
   x = 1e39; printf("x = 1e39    x = %e\n", x);

   x = 123.45123451234512345;
   printf("x = 123.45123451234512345     x = %3.17f\n", x);
}
\end{lstlisting}
\risultato
\codicepiccolo
\begin{console}
$ ./float
x = 1e-37    x = 1.000000e-37
x = 1e-38    x = 9.999999e-39
x = 1e-39    x = 1.000000e-39
x = 1e-40    x = 9.999946e-41
x = 1e-41    x = 9.999666e-42
x = 1e-42    x = 1.000527e-42
x = 1e-43    x = 9.949219e-44
x = 1e-44    x = 9.809089e-45
x = 1e-45    x = 1.401298e-45
x = 1e-46    x = 0.000000e+00
x = 1e38    x = 1.000000e+38
x = 1e39    x = inf
x = 123.45123451234512345
    x = 123.45123291015625000
\end{console}
\end{affianco}

(L'ultima riga dell'output è stata spezzata in due per farla stare nella
pagina.) Che cosa è successo? La figura~\ref{fig:retta-float} riassume i
valori positivi di un \code{float}:

\begin{itemize}
  \item fino a \code{FLT\_MIN}, circa $1.18 \times 10^{-38}$, i valori sono
        rappresentati con tutta la precisione: \code{1e-38} diventa
        \code{9.999999e-39}, cioè il \code{float} più vicino;
  \item sotto \code{FLT\_MIN} i numeri si possono ancora rappresentare, ma
        con \textbf{sempre meno cifre esatte} (sono detti
        \textit{subnormali}): \code{9.999946e-41}, \code{9.949219e-44},
        \code{9.809089e-45};
  \item il più piccolo \code{float} positivo è circa $1.4 \times 10^{-45}$:
        \code{1e-45} diventa \code{1.401298e-45}, e \code{1e-46}, ancora più
        piccolo, diventa \textbf{0};
  \item sopra \code{FLT\_MAX}, circa $3.4 \times 10^{38}$, non c'è più
        posto: \code{1e39} diventa \code{inf}, \textbf{infinito}.
\end{itemize}

\begin{figure}[H]
  \centering
  \begin{tikzpicture}[font=\Lato\scriptsize,
      tick/.style={font=\FiraCodeNL\tiny, anchor=north, align=center, inner sep=1pt}]
    \fill[red-gradient-1!40!white] (0,0) rectangle (1.4,0.55);
    \fill[orange-gradient-1!50!white] (1.4,0) rectangle (5.0,0.55);
    \fill[green-gradient-1!45!white] (5.0,0) rectangle (11.0,0.55);
    \fill[red-gradient-1!40!white] (11.0,0) rectangle (12.6,0.55);
    \draw[neutral!50!black, line width=0.6pt] (0,0) rectangle (12.6,0.55)
      (1.4,0) -- (1.4,0.55) (5.0,0) -- (5.0,0.55) (11.0,0) -- (11.0,0.55);
    \node at (0.7,0.275) {\textbf{0}};
    \node at (3.2,0.275) {\textbf{subnormali}};
    \node at (8.0,0.275) {\textbf{normali}};
    \node at (11.8,0.275) {\textbf{inf}};
    \node[tick] at (1.4,-0.05) {1.4e-45};
    \node[tick] at (5.0,-0.05) {1.18e-38\\(FLT\_MIN)};
    \node[tick] at (11.0,-0.05) {3.4e+38\\(FLT\_MAX)};
    \node[anchor=south, align=center] at (0.7,0.6) {troppo piccoli:\\diventano 0};
    \node[anchor=south, align=center] at (3.2,0.6) {sempre meno\\cifre esatte};
    \node[anchor=south, align=center] at (8.0,0.6) {circa 7 cifre esatte};
    \node[anchor=south, align=center] at (11.8,0.6) {troppo grandi:\\infinito};
  \end{tikzpicture}
  \caption{I valori positivi di un \code{float} (la scala non è
           proporzionale). Per i negativi vale lo stesso, con il segno meno.}
  \label{fig:retta-float}
\end{figure}

L'ultima riga mostra la \textbf{precisione}: del numero
\code{123.45123451234512345} il \code{float} conserva le prime cifre, poi
inventa. La figura~\ref{fig:cifre} confronta le cifre con quelle di un
\code{double}, che di cifre giuste ne conserva più del doppio.

\begin{figure}[H]
  \centering
  \begin{tikzpicture}[font=\FiraCodeNL\small,
      lab/.style={font=\Lato\scriptsize, anchor=east, text=neutral!35!black},
      giu/.style={anchor=base west, inner sep=0pt, text=mok, font=\FiraCodeNL\small\bfseries},
      sba/.style={anchor=base west, inner sep=0pt, text=merr, font=\FiraCodeNL\small}]
    \node[lab] at (-0.2,0.05) {nel programma};
    \node[anchor=base west, inner sep=0pt] at (0,0) {123.45123451234512345};
    \node[lab] at (-0.2,-0.6) {float};
    \node[giu] (f1) at (0,-0.65) {123.45123};
    \node[sba] at (f1.base east) {291015625000};
    \node[lab] at (-0.2,-1.25) {double};
    \node[giu] (d1) at (0,-1.3) {123.45123451234512};
    \node[sba] at (d1.base east) {402};
    \node[mnota, anchor=west] at (4.6,-0.6) {8 cifre giuste (ne sono garantite 6)};
    \node[mnota, anchor=west] at (4.6,-1.25) {17 cifre giuste (ne sono garantite 15)};
  \end{tikzpicture}
  \caption{Lo stesso numero in un \code{float} e in un \code{double}: in
           verde le cifre esatte, in rosso quelle sbagliate.}
  \label{fig:cifre}
\end{figure}

\subsection{Ai confini del \code{double}}

La stessa prova con un \code{double} (nell'ultima riga le cifre stampate sono
17, con \code{\%.17e}). Questa volta alcune costanti sono così estreme che
\code{gcc} protesta già in compilazione:

\begin{affianco}[0.56][t]
\codicepiccolo
\begin{lstlisting}[style=c, name=double.c]
#include <stdio.h>

int main()
{
   double x;

   x = 1e-307; printf("x = 1e-307    x = %e\n", x);
   x = 1e-308; printf("x = 1e-308    x = %e\n", x);
   x = 1e-309; printf("x = 1e-309    x = %e\n", x);
   x = 1e-400; printf("x = 1e-400    x = %e\n", x);
   x = 1e-401; printf("x = 1e-401    x = %e\n", x);
   x = 1e-402; printf("x = 1e-402    x = %e\n", x);
   x = 1e-403; printf("x = 1e-403    x = %e\n", x);
   x = 1e-404; printf("x = 1e-404    x = %e\n", x);
   x = 1e-405; printf("x = 1e-405    x = %e\n", x);
   x = 1e-406; printf("x = 1e-406    x = %e\n", x);

   x = 1e308; printf("x = 1e308    x = %e\n", x);
   x = 1e309; printf("x = 1e309    x = %e\n", x);

   x = 123.45123451234512345;
   printf("x = 123.45123451234512345     x = %.17e\n", x);
}
\end{lstlisting}
\risultato
\codicepiccolo
\begin{console}
$ ./double
x = 1e-307    x = 1.000000e-307
x = 1e-308    x = 1.000000e-308
x = 1e-309    x = 1.000000e-309
x = 1e-400    x = 0.000000e+00
x = 1e-401    x = 0.000000e+00
x = 1e-402    x = 0.000000e+00
x = 1e-403    x = 0.000000e+00
x = 1e-404    x = 0.000000e+00
x = 1e-405    x = 0.000000e+00
x = 1e-406    x = 0.000000e+00
x = 1e308    x = 1.000000e+308
x = 1e309    x = inf
x = 123.45123451234512345
    x = 1.23451234512345124e+02
\end{console}
\end{affianco}

Le costanti da \code{1e-400} a \code{1e-406} sono troppo piccole anche per
un \code{double}, e \code{1e309} è troppo grande: il compilatore lo dice con
un avviso per ciascuna riga (qui ne mostriamo due su otto) e le sostituisce
con 0 e con infinito.

\begin{console}
$ gcc -Wall -o double double.c
double.c: In function (*@‘@*)main(*@’@*):
double.c:10:4: warning: floating constant truncated to zero [-Woverflow]
   10 |    x = 1e-400; printf("x = 1e-400    x = %e\n", x);
      |    ^
double.c:19:4: warning: floating constant exceeds range of (*@‘@*)double(*@’@*) [-Woverflow]
   19 |    x = 1e309; printf("x = 1e309    x = %e\n", x);
      |    ^
\end{console}

\subsection{I reali non sono esatti}

Con un numero finito di bit si rappresenta un numero finito di valori: la
maggior parte dei numeri reali \textbf{non è rappresentabile}, e viene
sostituita dal valore rappresentabile più vicino. Succede anche a numeri
semplicissimi come $0.1$: in base 2 è un numero periodico
($0.000110011001100\dots$), esattamente come $1/3 = 0.333\dots$ in base 10,
e da qualche parte va troncato. Di conseguenza le operazioni con i reali,
a differenza di quelle con gli interi, non sono necessariamente esatte:

\begin{affianco}[0.5][c]
\begin{lstlisting}[style=c, name=decimali.c]
#include <stdio.h>

int main(void)
{
   double a = 0.1, b = 0.2;

   printf("0.1       = %.20f\n", a);
   printf("0.1 + 0.2 = %.20f\n", a + b);
   printf("0.3       = %.20f\n", 0.3);
   if (a + b == 0.3)
      printf("uguali\n");
   else
      printf("diversi!\n");
   return 0;
}
\end{lstlisting}
\risultato
\begin{console}
$ ./decimali
0.1       = 0.10000000000000000555
0.1 + 0.2 = 0.30000000000000004441
0.3       = 0.29999999999999998890
diversi!
\end{console}
\end{affianco}

$0.1 + 0.2$ e $0.3$ differiscono nella diciassettesima cifra, ma per
\code{==} sono diversi. Per questo due reali \textbf{non si confrontano con
\code{==}}: si controlla che la loro differenza sia piccola, per esempio con
\code{fabs(a + b - 0.3) < 1e-9} (\code{fabs}, il valore assoluto, è descritta
più avanti). Per i calcoli di piccola taglia questi errori sono trascurabili;
nei problemi grandi (la risoluzione di grandi sistemi di equazioni
differenziali, la minimizzazione di funzioni complicate) si accumulano, e
controllarli è il compito dell'analisi numerica.

\section{Dare un nome ai tipi: \code{typedef}}

Il meccanismo del \code{typedef} associa un identificatore a un tipo; da quel
momento l'identificatore si usa per dichiarare variabili (o funzioni)
esattamente come i tipi standard:

\begin{affianco}[0.5][c]
\begin{lstlisting}[style=c]
typedef char          uppercase;
typedef int           inches;
typedef unsigned long size_t;
typedef unsigned char u_char;

uppercase u;
u_char    c;
inches    length, width;
\end{lstlisting}
\risultato
A che cosa serve?
\begin{itemize}
  \item ad \textbf{abbreviare} le dichiarazioni lunghe (\code{u\_char} al
        posto di \code{unsigned char});
  \item a scrivere nomi che dicono \textbf{l'uso} delle variabili:
        \code{length} e \code{width} sono lunghezze in pollici
        (\textit{inches});
  \item a gestire le \textbf{differenze fra sistemi}.
\end{itemize}
\end{affianco}

Il terzo punto si capisce con un esempio. Un programma scritto per un sistema
con \code{int} di 4 byte va portato su un sistema dove gli \code{int} hanno
2 byte, troppo pochi per i suoi dati. Se il programma dichiara le sue
variabili con un tipo proprio, basta cambiare \textbf{una riga}:

\begin{affianco}[0.5][c]
\begin{lstlisting}[style=c]
/* sul primo sistema */
typedef int tipoInt;
\end{lstlisting}
\risultato
\begin{lstlisting}[style=c]
/* sul secondo sistema */
typedef long int tipoInt;
\end{lstlisting}
\end{affianco}

La libreria standard fa proprio così: \code{size\_t}, il tipo dei risultati di
\code{sizeof}, è definito con un \code{typedef} e cambia da un sistema
all'altro. Il \code{typedef} sarà molto comodo con le strutture
(capitolo~\ref{chap:strutture}).

\section{L'operatore \code{sizeof}}

\subsection{Quanti byte}

Il C fornisce l'operatore unario \code{sizeof}, che calcola il numero di byte
necessari per rappresentare un oggetto. L'oggetto, fra parentesi, può essere
un tipo (come \code{int} o \code{float}), un'espressione (come
\code{a + b}), un array o una struttura.
Ha la stessa priorità degli altri operatori unari. Il programma delle lezioni
lo applica ai tipi fondamentali:

\begin{lstlisting}[style=c, name=sizeofTypes.c]
/* Compute the size of some fundamental types. */

#include <stdio.h>

int main(void)
{
   printf("The size of some fundamental types is computed.\n\n");
   printf("       char:%3ld byte \n", sizeof(char));
   printf("      short:%3ld bytes\n", sizeof(short));
   printf("        int:%3ld bytes\n", sizeof(int));
   printf("       long:%3ld bytes\n", sizeof(long));
   printf("   unsigned:%3ld bytes\n", sizeof(unsigned));
   printf("      float:%3ld bytes\n", sizeof(float));
   printf("     double:%3ld bytes\n", sizeof(double));
   printf("long double:%3ld bytes\n", sizeof(long double));
   return 0;
}
\end{lstlisting}

\begin{console}
$ ./sizeofTypes
The size of some fundamental types is computed.

       char:  1 byte
      short:  2 bytes
        int:  4 bytes
       long:  8 bytes
   unsigned:  4 bytes
      float:  4 bytes
     double:  8 bytes
long double: 16 bytes
\end{console}

\code{sizeof(char)} vale 1 per definizione; gli altri valori sono quelli di
Linux su un PC a 64 bit e possono cambiare altrove. Il risultato di
\code{sizeof} è di tipo \code{size\_t}, un intero senza segno; la conversione
di \code{printf} pensata apposta è \code{\%zu}, ma su Linux a 64 bit
\code{size\_t} è un \code{unsigned long} e \code{\%ld} funziona.

\subsection{Una sorpresa}

Con \code{sizeof} verifichiamo quello che avevamo anticipato sulle costanti
carattere:

\begin{affianco}[0.56][c]
\begin{lstlisting}[style=c, name=dimensioni.c]
#include <stdio.h>

int main(void)
{
   char c = 'a';

   printf("sizeof(c)     = %zu\n", sizeof(c));
   printf("sizeof('a')   = %zu\n", sizeof('a'));
   printf("sizeof(c + c) = %zu\n", sizeof(c + c));
   return 0;
}
\end{lstlisting}
\risultato
\begin{console}
$ ./dimensioni
sizeof(c)     = 1
sizeof('a')   = 4
sizeof(c + c) = 4
\end{console}
\end{affianco}

Una variabile \code{char} occupa 1 byte, ma la costante \code{'a'} ne
occupa 4: è un \code{int}. E anche \code{c + c} è un \code{int}, per una
regola che vedremo alla fine del capitolo, la \textit{promozione a intero}.
Si noti infine che \code{sizeof} guarda solo il \textbf{tipo}
dell'espressione, senza calcolarla: \code{sizeof(i++)} non incrementa
\code{i}.

\section{Leggere e scrivere un carattere alla volta}

\subsection{\code{getchar} e \code{putchar}}

Per leggere e scrivere un carattere alla volta la libreria standard offre
due funzioni, dichiarate in \code{stdio.h}:

\begin{itemize}
  \item \code{getchar()} legge il prossimo carattere dallo
        \textbf{standard input} (la tastiera, o un file se l'input è
        rediretto con \code{<}) e ne restituisce il codice;
  \item \code{putchar(c)} scrive sullo \textbf{standard output} il
        carattere di codice \code{c}.
\end{itemize}

Di solito sono realizzate come macro (direttive per il preprocessore), ma si
usano come funzioni. E lavorano con gli \code{int}, perché un \code{int} può
rappresentare qualsiasi carattere e, in più, un valore speciale che segnala
la fine dell'input.

Il programma seguente legge una sequenza di caratteri e scrive ciascuno
\textbf{due volte}. Lo proviamo redirigendo sullo standard input il file
\code{input.txt}:

\begin{affianco}[0.5][c]
\begin{lstlisting}[style=c, name=double_out.c]
#include <stdio.h>

int main(void)
{
   int c;

   while ((c = getchar()) != EOF) {
      putchar(c);
      putchar(c);
   }
   return 0;
}
\end{lstlisting}
\risultato
\begin{console}
$ cat input.txt
abcdefghijklmno
$ ./double_out < input.txt
aabbccddeeffgghhiijjkkllmmnnoo

$
\end{console}
\end{affianco}

Da dove viene la riga vuota? Il file termina con un a capo, e anche
\code{'\textbackslash n'} è un carattere: raddoppiato, produce una riga
vuota. Se invece si lancia il programma senza redirezione, i caratteri
arrivano dalla tastiera:

\begin{console}
$ ./double_out
<<ciao>>
cciiaaoo

<<Ctrl+D>>
$
\end{console}

Il programma riceve i caratteri solo dopo \tasti{Invio}, tutti insieme: il
terminale raccoglie la riga (così la si può correggere) e la consegna quando
la si conferma. Poi \code{getchar} li restituisce uno alla volta. Per finire,
come nel capitolo~\ref{chap:primi-programmi}, si preme \tasti{Ctrl,D} a inizio
riga.

\subsection{Il ciclo di lettura}

La condizione \code{(c = getchar()) != EOF} fa tre cose in una
(figura~\ref{fig:ciclo-getchar}): legge un carattere, lo assegna a \code{c} e
controlla che l'input non sia finito.

\begin{figure}[H]
  \centering
  \setlength{\fboxsep}{1.5pt}
  \renewcommand{\arraystretch}{1.6}
  \begin{tabular}{@{}c>{\FiraCodeNL\small}l>{\raggedright\arraybackslash}p{8.4cm}@{}}
    \wtpassonum{1} & while ((c = \colorbox{accent!16!white}{getchar()}) != EOF)
                   & \code{getchar()} legge un carattere, per esempio \code{a}, e restituisce il suo codice, 97 \\
    \wtpassonum{2} & while ((\colorbox{accent!16!white}{c = 97}) != EOF)
                   & il valore viene assegnato a \code{c}; l'assegnazione vale 97 \\
    \wtpassonum{3} & while (\colorbox{accent!16!white}{97 != EOF})
                   & 97 è diverso da \code{EOF}: la condizione è vera \\
    \wtpassonum{4} & while (1)
                   & si esegue il corpo, poi si ricomincia dal passo 1 \\
  \end{tabular}
  \caption{Il calcolo della condizione, passo per passo. Quando l'input
           finisce, \code{getchar()} restituisce \code{EOF}, il confronto
           \code{EOF != EOF} è falso e il ciclo termina.}
  \label{fig:ciclo-getchar}
\end{figure}

Tutte le parentesi sono necessarie. L'operatore \code{!=} ha la precedenza
su \code{=} (capitolo~\ref{chap:sintassi}), quindi senza le parentesi interne
\code{c = getchar() != EOF} verrebbe letto come
\code{c = (getchar() != EOF)}: a \code{c} andrebbe il risultato del
confronto, 1 o 0, e il programma stamperebbe, al posto di ogni carattere,
due caratteri di codice 1, invisibili. \code{gcc} con \code{-Wall} intuisce
l'errore e suggerisce le parentesi:

\begin{console}
$ gcc -Wall -o senza_parentesi senza_parentesi.c
senza_parentesi.c: In function (*@‘@*)main(*@’@*):
senza_parentesi.c:7:11: warning: suggest parentheses around assignment used as truth value [-Wparentheses]
    7 |    while (c = getchar() != EOF) {
      |           ^
\end{console}

\subsection{\code{EOF}}

\code{EOF} (\textit{end of file}, fine del file) è il valore che
\code{getchar} restituisce quando l'input è finito: alla fine di un file
rediretto, o quando si preme \tasti{Ctrl,D} a inizio riga. È una costante
simbolica definita in \code{stdio.h}, tipicamente come

\begin{lstlisting}[style=c]
#define EOF (-1)
\end{lstlisting}

ma su un altro sistema potrebbe valere diversamente: usare il nome
\code{EOF} invece di $-1$ garantisce la \textbf{portabilità}. Stampandolo:

\begin{affianco}[0.5][c]
\begin{lstlisting}[style=c, name=eof.c]
#include <stdio.h>

int main()
{
   printf("EOF = %d\n", EOF);
   printf("EOF = %c\n", EOF);
}
\end{lstlisting}
\risultato
\begin{console}
$ ./eof
EOF = -1
EOF = (*@\char"FFFD@*)
\end{console}
\end{affianco}

Con \code{\%d} si vede $-1$. Con \code{\%c} viene scritto un byte che vale
255 (\code{11111111}, la lettura senza segno di $-1$), che non corrisponde a
nessun carattere: il terminale mostra un simbolo di carattere non valido.

\subsection{Perché \code{int} e non \code{char}}

Nel programma, \code{c} è dichiarata \code{int} e non \code{char}. Il motivo
è che \code{getchar} può restituire \textbf{257 valori} diversi: i 256
possibili byte, da 0 a 255, più \code{EOF}. Un \code{char} ne contiene solo
256, e due valori finirebbero per confondersi (figura~\ref{fig:getchar-int}):
\code{char} può essere \code{signed} o \code{unsigned}, e \code{EOF} potrebbe
non valere $-1$; un \code{int} va bene in ogni caso.

\begin{figure}[H]
  \centering
  \begin{tikzpicture}[font=\FiraCodeNL\scriptsize,
      gc/.style={mcella, minimum width=0.95cm, minimum height=0.55cm}]
    \noindici
    \celle[gc]{g}{(0,0)}{EOF,0,1,2,\ldots,254,255}
    \node[mindice, anchor=north] at ([yshift=-1pt]g-0.south) {$-1$};
    \draw[mrif, draw=merr] (g-6.north) to[out=140, in=40]
      node[above, mnota, text=merr] {in un \code{signed char} il byte 255 diventa $-1$: sembra \code{EOF}} (g-0.north);
    \draw[mrif, draw=merr] ([yshift=-0.3cm]g-0.south) to[out=-40, in=-140]
      node[below, mnota, text=merr] {in un \code{unsigned char} $-1$ diventa 255: \code{c} non vale mai \code{EOF}} ([yshift=-0.1cm]g-6.south);
    \node[mnota, anchor=west, align=left] at ([xshift=0.5cm]g-6.east)
      {257 valori possibili:\\in un \code{int} restano tutti distinti \mok};
  \end{tikzpicture}
  \caption{I valori che \code{getchar} può restituire, e che cosa succede se
           li si mette in un \code{char}.}
  \label{fig:getchar-int}
\end{figure}

Con \code{unsigned char} il guasto è evidente:

\begin{affianco}[0.5][c]
\begin{lstlisting}[style=c, name=uchar.c]
#include <stdio.h>

int main(void)
{
   unsigned char c;

   while ((c = getchar()) != EOF) {
      putchar(c);
      putchar(c);
   }
   return 0;
}
\end{lstlisting}
\risultato
\begin{console}
$ ./uchar < input.txt
aabbccddeeffgghhiijjkkllmmnnoo

(*@\char"FFFD\char"FFFD\char"FFFD\char"FFFD\char"FFFD\char"FFFD\char"FFFD\char"FFFD\char"FFFD\char"FFFD\char"FFFD\char"FFFD\char"FFFD\char"FFFD\char"FFFD\char"FFFD\char"FFFD\char"FFFD\char"FFFD\char"FFFD@*)
^C
\end{console}
\end{affianco}

Dopo il contenuto del file il programma scrive all'infinito il byte 255
(mostrato come carattere non valido), finché non lo si ferma con
\tasti{Ctrl,C}. Il motivo: quando l'input finisce, \code{getchar}
restituisce $-1$, che in \code{c} diventa 255; il confronto
\code{255 != EOF} è sempre vero, e il ciclo non termina più. Il compilatore
lo segnala solo con l'opzione \code{-Wextra}, che attiva altri avvisi oltre a
quelli di \code{-Wall}:

\begin{console}
$ gcc -Wall -Wextra -o uchar uchar.c
uchar.c: In function (*@‘@*)main(*@’@*):
uchar.c:7:27: warning: comparison is always true due to limited range of data type [-Wtype-limits]
    7 |    while ((c = getchar()) != EOF) {
      |                           ^~
\end{console}

Con un \code{signed char} il guasto è più subdolo: il programma funziona,
salvo fermarsi in anticipo se l'input contiene un byte che vale 255.

\subsection{Maiuscole e minuscole}

I caratteri sono ordinati come gli interi che li rappresentano, e si
comportano come tali:

\begin{itemize}
  \item \code{'a' + 1} vale \code{'b'}, e \code{'b' + 1} vale \code{'c'};
  \item \code{'z' - 'a'} vale 25;
  \item \code{'A' < 'a'}, perché 65 è minore di 97;
  \item \code{'A' - 'a'} ha lo stesso valore di \code{'B' - 'b'}: la
        distanza fra maiuscola e minuscola è la stessa per tutte le lettere
        (figura~\ref{fig:maiuscole}).
\end{itemize}

\begin{figure}[H]
  \centering
  \begin{tikzpicture}[font=\FiraCodeNL\scriptsize,
      lc/.style={mcella, minimum width=0.8cm, minimum height=0.55cm, font=\FiraCodeNL\small\bfseries},
      cod/.style={font=\FiraCodeNL\tiny, text=neutral!45!black, inner sep=1pt}]
    \noindici
    \foreach \mi/\ma/\nmi/\nma [count=\k from 0] in {a/A/97/65, b/B/98/66, c/C/99/67,
        d/D/100/68, e/E/101/69, f/F/102/70, g/G/103/71, \ldots/\ldots/{}/{}, y/Y/121/89, z/Z/122/90} {
      \node[lc, fill=accent!12!white, draw=accent!60!black] (m\k) at (\k*1.15,0) {\mi};
      \node[cod, anchor=south] at (m\k.north) {\nmi};
      \node[lc, fill=primary!12!white] (M\k) at (\k*1.15,-1.5) {\ma};
      \node[cod, anchor=north] at (M\k.south) {\nma};
    }
    \foreach \k in {0,1,2,3,4,5,6,8,9}
      \draw[wtarrowacc] (m\k.south) -- (M\k.north);
    \node[mnota, anchor=west, align=left] at ([xshift=0.4cm]$(m9.east)!0.5!(M9.east)$)
      {ogni freccia somma\\\code{'A' - 'a'}, cioè $-32$};
  \end{tikzpicture}
  \caption{Minuscole e maiuscole (sotto ciascuna, il codice): ogni lettera
           dista 32 posizioni dalla sua maiuscola.}
  \label{fig:maiuscole}
\end{figure}

Su questa regolarità si basa \code{capitalize.c}, che trasforma in maiuscolo
le lettere minuscole e copia senza modifiche tutti gli altri caratteri:

\begin{affianco}[0.5][c]
\begin{lstlisting}[style=c, name=capitalize.c]
#include <stdio.h>

int main(void)
{
   int   c;

   while ((c = getchar()) != EOF)
      if (c >= 'a' && c <= 'z')
         putchar(c + 'A' - 'a');
      else
         putchar(c);
   return 0;
}
\end{lstlisting}
\risultato
\begin{console}
$ cat input.txt
abcDE6ygTr56
$ ./capitalize < input.txt
ABCDE6YGTR56
\end{console}
\end{affianco}

Per la \code{g}, per esempio: $103 + 65 - 97 = 71$, il codice di
\code{'G'}. Scrivere \code{'A' - 'a'} invece di $-32$ rende il programma più
leggibile e non obbliga a ricordare la tabella. Il corpo del \code{while} è
un'unica istruzione, l'\code{if}-\code{else}, quindi le graffe non servono.
(La libreria standard ha una funzione che fa la stessa cosa, \code{toupper},
dichiarata in \code{ctype.h}.)

\section{Funzioni matematiche}

\subsection{La libreria matematica}

Il linguaggio C non contiene funzioni matematiche predefinite: le funzioni
matematiche si trovano nella \textbf{libreria matematica}, che fa parte della
libreria standard. Per esempio \code{sqrt} (radice quadrata), \code{pow}
(potenza), \code{log}, \code{sin} e \code{cos}. In genere ricevono parametri
di tipo \code{double} e restituiscono un risultato \code{double}; i
prototipi sono in \code{math.h}, e su Linux la libreria va collegata con
l'opzione \code{-lm} (capitolo~\ref{chap:primi-programmi}). Il programma
seguente usa \code{sqrt}, che ha un parametro, e \code{pow}, che ne ha due:

\begin{affianco}[0.56][t]
\codicepiccolo
\begin{lstlisting}[style=c, name=radice.c]
#include <math.h>
#include <stdio.h>

int main(void)
{
   double x;

   printf("\n%s\n%s\n%s\n\n",
      "The square root of x and x raised",
      "to the x power will be computed.",
      "---");
   while (1) {                  /* do it forever */
      printf("Input x:  ");
      scanf("%lf", &x);
      if (x >= 0.0)
         printf("\n%15s%22.15e\n%15s%22.15e\n%15s%22.15e\n\n",
            "x = ", x,
            "sqrt(x) = ", sqrt(x),
            "pow(x, x) = ", pow(x, x));
      else
         printf("\nSorry, your number must be nonnegative.\n\n");
   }
   return 0;
}
\end{lstlisting}
\risultato
\begin{console}
$ gcc -Wall -o radice radice.c -lm
$ ./radice

The square root of x and x raised
to the x power will be computed.
---

Input x:  <<2>>

           x =  2.000000000000000e+00
     sqrt(x) =  1.414213562373095e+00
   pow(x, x) =  4.000000000000000e+00

Input x:  <<-1>>

Sorry, your number must be nonnegative.

Input x:  ^C
\end{console}
\end{affianco}

\begin{itemize}
  \item La prima \code{printf} stampa tre stringhe passate come argomenti, con
        la conversione \code{\%s}: è un modo per spezzare un testo lungo.
  \item \code{\%15s} scrive una stringa in 15 caratteri, allineata a destra;
        \code{\%22.15e} scrive un numero in notazione scientifica con 15
        cifre dopo il punto, in 22 caratteri (capitolo~\ref{chap:primi-programmi}).
        Per questo i risultati sono incolonnati.
  \item \code{while (1)} non termina mai: si esce con \tasti{Ctrl,C}. Qui
        \tasti{Ctrl,D} non basta: \code{scanf} non legge più nulla,
        \code{x} conserva l'ultimo valore e il ciclo si ripete all'infinito.
  \item Nelle slide l'esponente è scritto con tre cifre, \code{e+000}: è il
        formato usato su Windows; su Linux le cifre sono due.
\end{itemize}

\subsection{Il valore assoluto}

Per il valore assoluto la libreria standard ha due funzioni:
\code{abs}, per gli \code{int}, con il prototipo in \code{stdlib.h}, e
\code{fabs}, per i \code{double}, con il prototipo in \code{math.h}.
Scambiarle è un errore insidioso:

\begin{affianco}[0.5][c]
\begin{lstlisting}[style=c, name=assoluto.c]
#include <stdio.h>
#include <stdlib.h>
#include <math.h>

int main(void)
{
   printf("abs(-7)    = %d\n", abs(-7));
   printf("fabs(-2.7) = %f\n", fabs(-2.7));
   printf("abs(-2.7)  = %d\n", abs(-2.7));
   return 0;
}
\end{lstlisting}
\risultato
\begin{console}
$ ./assoluto
abs(-7)    = 7
fabs(-2.7) = 2.700000
abs(-2.7)  = 2
\end{console}
\end{affianco}

\code{abs(-2.7)} vale 2: il prototipo di \code{abs} dice che il parametro è
un \code{int}, quindi $-2.7$ viene convertito in intero ($-2$) prima della
chiamata. Con \code{-Wall} non c'è alcun avviso; con \code{-Wextra} sì:

\begin{console}
$ gcc -Wall -Wextra -o assoluto assoluto.c -lm
assoluto.c: In function (*@‘@*)main(*@’@*):
assoluto.c:9:32: warning: using integer absolute value function (*@‘@*)abs(*@’@*) when argument is of floating-point type (*@‘@*)double(*@’@*) [-Wabsolute-value]
    9 |    printf("abs(-2.7)  = %d\n", abs(-2.7));
      |                                ^~~
\end{console}

\section{Conversioni e cast}

Un'espressione aritmetica come \code{x + y} ha un valore e un tipo. Se
\code{x} e \code{y} sono entrambi \code{int}, anche \code{x + y} è un
\code{int}. Ma che succede se i tipi sono diversi, o più piccoli di un
\code{int}? Il C applica delle \textbf{conversioni implicite}, secondo regole
precise; e il programmatore può chiederne di \textbf{esplicite}, con il
\textit{cast}.

\subsection{La promozione a intero}

Un \code{char} o uno \code{short}, con o senza segno, si possono usare in
qualunque espressione in cui si può usare un \code{int} o un
\code{unsigned int}. Nelle espressioni vengono infatti \textbf{promossi}:
se tutti i valori del tipo originale sono rappresentabili da un \code{int}
(ed è il caso dei PC), il valore viene convertito in \code{int}; altrimenti
in \code{unsigned int}. Il processore, del resto, lavora in modo naturale su
valori della dimensione di un \code{int}.

Così, se \code{x} e \code{y} sono \code{short}, \code{x + y} è un
\code{int}; e in

\begin{lstlisting}[style=c]
char c = 'A';
printf("%c\n", c);
\end{lstlisting}

il parametro passato a \code{printf} è un'espressione, \code{c}, che per la
promozione a intero è di tipo \code{int}, con valore 65. Lo stesso motivo
rende \code{sizeof(c + c)} uguale a 4.

\subsection{Le conversioni aritmetiche}

Quando un operatore binario (\code{+}, \code{-}, \code{*}, \code{/}, un
confronto\dots) ha operandi di tipo diverso, prima del calcolo uno dei due
viene convertito nel tipo dell'altro. Con \code{int i;} e \code{float f;},
in \code{i + f} il valore di \code{i} viene convertito in \code{float}, e il
risultato è un \code{float}. La regola generale è che l'operando di tipo
\textbf{“più piccolo”} viene convertito nel tipo \textbf{“più grande”}
(figura~\ref{fig:scala}). Nell'ordine:

\begin{enumerate}
  \item se uno dei due operandi è \code{long double}, l'altro viene
        convertito in \code{long double};
  \item altrimenti, se uno è \code{double}, l'altro diventa \code{double};
  \item altrimenti, se uno è \code{float}, l'altro diventa \code{float};
  \item altrimenti si applica la promozione a intero a entrambi, e poi:
        \begin{itemize}
          \item se uno è \code{unsigned long}, l'altro diventa
                \code{unsigned long};
          \item altrimenti, se uno è \code{long} e l'altro
                \code{unsigned}: se un \code{long} può rappresentare tutti i
                valori di un \code{unsigned}, l'\code{unsigned} diventa
                \code{long}, altrimenti entrambi diventano
                \code{unsigned long};
          \item altrimenti, se uno è \code{long}, l'altro diventa
                \code{long};
          \item altrimenti, se uno è \code{unsigned}, l'altro diventa
                \code{unsigned};
          \item altrimenti entrambi sono \code{int}.
        \end{itemize}
\end{enumerate}

\begin{figure}[H]
  \centering
  \begin{tikzpicture}[font=\FiraCodeNL\scriptsize,
      gr/.style={draw=primary!60!black, fill=primary!6!white, rounded corners=2pt,
        line width=0.7pt, minimum width=2.7cm, minimum height=0.5cm},
      grr/.style={gr, draw=accent!60!black, fill=accent!8!white}]
    \foreach \testo [count=\k from 0] in {int, unsigned, long, unsigned long}
      \node[gr] (s\k) at (0,\k*0.8) {\testo};
    \foreach \testo [count=\k from 4] in {float, double, long double}
      \node[grr] (s\k) at (0,\k*0.8) {\testo};
    \foreach \k [evaluate=\k as \kk using int(\k+1)] in {0,...,5}
      \draw[wtarrow] (s\k.north) -- (s\kk.south);
    \node[wtnodeplain, font=\FiraCodeNL\scriptsize] (cs) at (0,-1.15) {char, short};
    \draw[wtarrow, draw=accent!70!black] (cs.north) -- node[right, mnota] {promozione a intero} (s0.south);
    \draw[wtbrace] ([xshift=0.3cm]s0.south east) -- ([xshift=0.3cm]s3.north east);
    \node[mnota, anchor=west] at ([xshift=0.55cm]$(s0.east)!0.5!(s3.east)$) {tipi interi};
    \draw[wtbrace] ([xshift=0.3cm]s4.south east) -- ([xshift=0.3cm]s6.north east);
    \node[mnota, anchor=west] at ([xshift=0.55cm]s5.east) {tipi reali};
    \node[mnota, anchor=east, align=right] at ([xshift=-0.5cm]s3.west)
      {fra due operandi, quello più in basso\\viene convertito nel tipo dell'altro};
    \node[mnota, anchor=east, align=right, text=neutral!45!black] at ([xshift=-0.5cm]$(s1.west)!0.5!(s2.west)$)
      {\code{unsigned} con \code{long}: \code{long} su Linux e macOS\\a 64 bit, \code{unsigned long} su Windows};
  \end{tikzpicture}
  \caption{La scala delle conversioni aritmetiche (come su Linux a 64 bit).}
  \label{fig:scala}
\end{figure}

Ecco l'esempio delle slide. Con le dichiarazioni

\begin{lstlisting}[style=c]
char c;  short s;  int i;  long l;  unsigned u;  unsigned long ul;
float f;  double d;  long double ld;
\end{lstlisting}

le espressioni seguenti hanno questi tipi (verificati con il compilatore,
su Linux a 64 bit):

\begin{center}
\renewcommand{\arraystretch}{1.12}
\begin{tabular}{@{}>{\FiraCodeNL}l l@{\hspace{1.4cm}}>{\FiraCodeNL}l l@{}}
  \toprule
  \normalfont\textbf{espressione} & \textbf{tipo} & \normalfont\textbf{espressione} & \textbf{tipo} \\
  \midrule
  c - s / i   & \code{int}    & u * 7 - i   & \code{unsigned} \\
  u * 2.0 + 1 & \code{double} & f * 7 - i   & \code{float} \\
  c + 3       & \code{int}    & 7 * s * ul  & \code{unsigned long} \\
  c + 5.0     & \code{double} & ld + c      & \code{long double} \\
  d + s       & \code{double} & u - ul      & \code{unsigned long} \\
  2 * i / l   & \code{long}   & u - l       & dipende dal sistema \\
  \bottomrule
\end{tabular}
\end{center}

Qualche caso commentato: in \code{c - s / i}, \code{c} e \code{s} vengono
promossi a \code{int} e tutto il calcolo è fra \code{int}; in
\code{u * 7 - i} la costante 7 e poi \code{i} vengono convertiti in
\code{unsigned}; in \code{u * 2.0 + 1} la costante \code{2.0} è un
\code{double}, e trascina tutto in \code{double}. L'ultimo dipende davvero dal
sistema: su Linux e macOS a 64 bit un \code{long} (8 byte) contiene tutti
i valori di un \code{unsigned} (4 byte), quindi \code{u - l} è un
\code{long}; su Windows, dove \code{long} ha 4 byte, è un
\code{unsigned long}.

\error{Interi con e senza segno}{
  Mescolare \code{int} e \code{unsigned} può dare risultati sorprendenti.
  In \code{-1 < 1u}, per la regola vista, $-1$ viene convertito in
  \code{unsigned} e diventa $4\,294\,967\,295$: il confronto è
  \textbf{falso}. Un programma con \code{int i = -1; unsigned u = 1;} e
  \code{if (i < u)} sceglie quindi il ramo \code{else}. Anche qui l'avviso
  arriva solo con \code{-Wextra}: \textit{comparison of integer expressions
  of different signedness}. Conviene compilare sempre con
  \code{-Wall -Wextra}.}

\subsection{Le conversioni nelle assegnazioni}

Anche l'assegnazione converte: il valore viene trasformato nel tipo della
variabile che lo riceve. Da reale a intero la parte decimale viene
\textbf{troncata}, non arrotondata (cioè si va verso lo zero); da intero a
reale non si perde nulla.

\begin{affianco}[0.5][c]
\begin{lstlisting}[style=c, name=tronca.c]
#include <stdio.h>

int main(void)
{
   int i = 3.9;
   int j = -3.9;
   double x = 7;

   printf("i = %d   j = %d   x = %f\n", i, j, x);
   return 0;
}
\end{lstlisting}
\risultato
\begin{console}
$ ./tronca
i = 3   j = -3   x = 7.000000
\end{console}
\end{affianco}

\subsection{Il cast}

Le conversioni viste finora sono \textbf{implicite}: le fa il compilatore da
solo. Esistono anche le conversioni \textbf{esplicite}, chiamate
\term{cast}: si scrive il tipo desiderato fra parentesi davanti a
un'espressione.

\begin{lstlisting}[style=c]
int i;
double f;

f = (double) i;
\end{lstlisting}

La variabile \code{i} \textbf{non viene modificata}: cambia il tipo del
valore prodotto dall'espressione \code{i}. L'uso tipico è evitare la
divisione intera, per esempio nel calcolo di una media:

\begin{affianco}[0.56][c]
\begin{lstlisting}[style=c, name=media.c]
#include <stdio.h>

int main(void)
{
   int somma = 7, n = 2;
   double media;

   media = somma / n;
   printf("somma / n            = %f\n", media);
   media = (double) somma / n;
   printf("(double) somma / n   = %f\n", media);
   media = (double) (somma / n);
   printf("(double) (somma / n) = %f\n", media);
   printf("somma vale ancora %d\n", somma);
   return 0;
}
\end{lstlisting}
\risultato
\begin{console}
$ ./media
somma / n            = 3.000000
(double) somma / n   = 3.500000
(double) (somma / n) = 3.000000
somma vale ancora 7
\end{console}
\end{affianco}

Nel primo caso la divisione è fra due \code{int} e vale 3; il risultato viene
convertito in \code{double} solo dopo, nell'assegnazione. Il cast è un
operatore unario e ha la precedenza su \code{/}: in
\code{(double) somma / n} si converte prima \code{somma} in $7.0$, poi la
divisione è fra un \code{double} e un \code{int}, quindi anche \code{n} viene
convertito e il risultato è $3.5$. Nel terzo caso le parentesi fanno
eseguire prima la divisione intera: il cast arriva quando il danno è fatto
(figura~\ref{fig:cast}).

\begin{figure}[H]
  \centering
  \begin{tikzpicture}[font=\FiraCodeNL\small,
      op/.style={circle, draw=accent!65!black, fill=accent!8!white, line width=0.8pt,
        minimum size=0.62cm, inner sep=0pt, font=\FiraCodeNL\small\bfseries},
      ca/.style={draw=accent!65!black, fill=accent!8!white, line width=0.8pt,
        rounded corners=0.3cm, minimum height=0.62cm, inner xsep=5pt, font=\FiraCodeNL\scriptsize\bfseries},
      fo/.style={wtnodeplain, font=\FiraCodeNL\scriptsize, minimum width=1.2cm},
      val/.style={font=\Lato\scriptsize\bfseries, text=green-gradient-6!50!black, anchor=west},
      ord/.style={circle, fill=primary!75!black, text=white, font=\Lato\tiny\bfseries,
        inner sep=0pt, minimum size=0.36cm},
      tit/.style={font=\FiraCodeNL\small}]
    % (double) somma / n
    \node[tit] at (0.4,1.0) {(double) somma / n};
    \node[op] (d1) at (0.4,0) {/};
    \node[ca] (c1) at (-0.9,-1.1) {(double)};
    \node[fo] (s1) at (-0.9,-2.2) {somma\\[-1pt]{\color{neutral!50!black}(7)}};
    \node[fo] (n1) at (1.7,-1.1) {n\\[-1pt]{\color{neutral!50!black}(2)}};
    \foreach \a/\b in {d1/c1, d1/n1, c1/s1} \draw[fsramo] (\a) -- (\b);
    \node[val] at ([xshift=4pt]c1.east) {= 7.0};
    \node[val] at ([xshift=4pt]d1.east) {= 7.0 / 2.0 = 3.5};
    \node[ord] at ([shift={(-0.25,0.3)}]c1.north west) {1};
    \node[ord] at ([shift={(-0.42,0.32)}]d1.center) {2};
    % (double) (somma / n)
    \node[tit] at (8.0,1.0) {(double) (somma / n)};
    \node[ca] (c2) at (8.0,0) {(double)};
    \node[op] (d2) at (8.0,-1.1) {/};
    \node[fo] (s2) at (6.8,-2.2) {somma\\[-1pt]{\color{neutral!50!black}(7)}};
    \node[fo] (n2) at (9.2,-2.2) {n\\[-1pt]{\color{neutral!50!black}(2)}};
    \foreach \a/\b in {c2/d2, d2/s2, d2/n2} \draw[fsramo] (\a) -- (\b);
    \node[val] at ([xshift=4pt]d2.east) {= 7 / 2 = 3 (divisione intera)};
    \node[val] at ([xshift=4pt]c2.east) {= 3.0};
    \node[ord] at ([shift={(-0.42,0.32)}]d2.center) {1};
    \node[ord] at ([shift={(-0.25,0.3)}]c2.north west) {2};
  \end{tikzpicture}
  \caption{Dove sta il cast cambia tutto: a sinistra converte \code{somma}
           prima della divisione, a destra il risultato già troncato.}
  \label{fig:cast}
\end{figure}

\insintesi{
  Il C ha dodici tipi fondamentali: i tipi \textbf{interi} (\code{char},
  \code{short}, \code{int}, \code{long}, con e senza segno) e i tipi
  \textbf{reali} (\code{float}, \code{double}, \code{long double}). Un
  \code{char} è un intero di 1 byte che contiene il \textbf{codice} di un
  carattere (ASCII: le cifre, le maiuscole e le minuscole hanno codici
  consecutivi); le \textbf{sequenze di escape} scrivono i caratteri
  speciali. Ogni tipo ha un intervallo di valori (\code{limits.h},
  \code{float.h}): superarlo con gli interi provoca un \textbf{overflow}
  silenzioso, i reali hanno una \textbf{precisione} limitata e non sono
  esatti, quindi non si confrontano con \code{==}. \code{sizeof} dà i byte
  di un tipo o di un'espressione, \code{typedef} dà un nuovo nome a un tipo.
  \code{getchar} e \code{putchar} leggono e scrivono un carattere; la
  variabile che riceve \code{getchar} deve essere \code{int}, per
  distinguere \code{EOF}. Le funzioni matematiche sono in \code{math.h}
  (\code{-lm}). Nelle espressioni \code{char} e \code{short} diventano
  \code{int}, e l'operando di tipo più piccolo viene convertito nel tipo
  dell'altro; il \textbf{cast} \code{(tipo)} converte esplicitamente.}

\section{Riferimenti}

\begin{itemize}
  \item \textbf{A. Kelley, I. Pohl}, \textit{A Book on C}: capitolo~3 (i
        tipi di dato fondamentali).
  \item \textbf{Arithmetic types} --- cppreference.com. \\
        \urlx{https://en.cppreference.com/w/c/language/arithmetic_types}
  \item \textbf{Implicit conversions} --- cppreference.com. \\
        \urlx{https://en.cppreference.com/w/c/language/conversion}
  \item \textbf{D. Goldberg}, \textit{What Every Computer Scientist Should
        Know About Floating-Point Arithmetic} (il classico sugli errori dei
        reali). \\
        \urlx{https://docs.oracle.com/cd/E19957-01/806-3568/ncg_goldberg.html}
\end{itemize}
